
\chapter{概形的态射}
\label{ch:morphisms-of-schemes}

我们现在来描述\trScheme{}之间的\trMorphism。我们将定义\trMorphism{}的一些易于陈述的性质，而把更微妙的性质留待以后。

回想：\trScheme{}是 (i) 一个\trSet，(ii) 带一个\trTopology，(iii) 加一个（\trStructureSheaf{}意义的）函数环层——定义可作三步观。同样，\trScheme{}的\trMorphism{} $X\to Y$ 亦可定义为：(i) \trSet{}的映射，(ii) 它连续，(iii) 再带上涉及函数层的进一步信息。\trAffineScheme{}的集合映射已见 \S3.2.10，其连续性已见\cref{ex:ring-maps-continuous-maps}。


\section{「正确」态射定义的动机}

这里给出两个动机，说明态射应当有怎样的行为。第一个是代数的，第二个是几何的。

\textbf{7.1.1.}\ \emph{代数动机。}\quad 我们将会希望：诸\trAffineScheme{}的%
\trMorphism{} $\Spec A\to\Spec B$ 恰好就是诸\trRingMorphism{} $B\to A$。我们已经%
见过，\trRingMorphism{} $B\to A$ 会朝相反的方向诱导诸\trTopologicalSpace{}之间的%
映射（\cref{ex:ring-maps-continuous-maps}）；主要的新配料，则是要掺进这份混合物%
里的函数\trStructureSheaf。于是，概形之间的\trMorphism{}就应当是这样一个东西——%
「在\trAffine\trOpenSet{}的层面上，看上去便是这副模样。」

\textbf{7.1.2.}\ \emph{几何动机。}\quad 受\trManifold{}理论（\S3.1.1）的驱动——%
\trManifold{}与\trScheme{}一样，都是\trRingedSpace——我们希望概形之间的%
\trMorphism{}至少得是\trRingedSpace{}之间的\trMorphism；至于这类态射为何物，我们%
现在就来为它寻觅动机。（我们将在下一节把它形式化。）注意：若 $\pi\colon X\to Y$ 是%
诸\trDifferentiable\trManifold{}之间的一个映射，则 $Y$ 上的一个光滑函数拉回到
$X$ 上，得到的仍是 $X$ 上的一个光滑函数。更精确地说，给定一个\trOpenSubset{}
$\mathsf{U}\subset Y$，就有一个自然的映射
$\Gamma(\mathsf{U},\OO_Y)\to\Gamma(\pi^{-1}(\mathsf{U}),\OO_X)$。这条映射与限制%
良好相容（把一个函数限制到更小的\trOpenSet{}上再拉回，与先拉回再限制，结果相同），
所以事实上我们有 $Y$ 上诸\trSheaf{}的一个映射 $\OO_Y\to\pi_*\OO_X$。同样地，概形%
之间的一个\trMorphism{} $\pi\colon X\to Y$ 也应当诱导一个映射
$\OO_Y\to\pi_*\OO_X$。但事实上，在诸\trDifferentiable\trManifold{}的\trCategory{}%
中，一个\trContinuousMap{} $\pi\colon X\to Y$ 是诸\trDifferentiable\trManifold{}%
之间的一个映射，当且仅当 $Y$ 上的光滑函数拉回到 $X$ 上仍是光滑函数（也就是说，从
$Y$ 上的光滑函数到 $X$ 上的\emph{函数}的那个\trPullback{}映射，实际上正落在%
\emph{光滑函数}这个子集之中；也就是说，\trContinuousMap{} $\pi$ 诱导了光滑函数的%
一个\trPullback，而它可以解读为一个映射 $\OO_Y\to\pi_*\OO_X$）。于是，这个%
\trSheaf{}的映射便\emph{刻画}了\trDifferentiable{}范畴中的态射。所以，我们本可以%
拿它当\trDifferentiable{}范畴中\trMorphism{}的\emph{定义}来用（见%
\cref{ex:pullback-detects-smoothness}）。

但我们要怎样把这套道理用到概形的\trCategory{}上去呢？在诸\trDifferentiable%
\trManifold{}的\trCategory{}中，一个\trContinuousMap{} $\pi\colon X\to Y$ 会%
\emph{诱导}函数之\trSheaf{}的一个\trPullback，我们可以问：它何时会诱导\emph{光滑}%
函数的一个\trPullback。然而，函数在概形上要更为古怪，单凭诸\trTopologicalSpace{}%
之间的映射，我们复原不出那个\trPullback{}映射。正确的补丁是把这一点硬性焊进%
\trMorphism{}的定义之中：也就是说，取一个\trContinuousMap{} $\pi\colon X\to Y$，
连同一条\trPullback{}映射 $\pi^{\sharp}\colon\OO_Y\to\pi_*\OO_X$。由此便得到赋环%
空间的\emph{\trCategory}的定义。

人们或许指望把概形之间的\trMorphism{}定义为\trRingedSpace{}之间的\trMorphism。
但这并不全对，因为那样一来，动机~7.1.1就不满足了：如我们所愿，对每个\trMorphism{}
$A\to B$ 都有\trMorphism{} $\Spec B\to\Spec A$，可是还会冒出额外的、并非这样产生
的\trRingedSpace{}的\trMorphism{} $\Spec B\to\Spec A$（见\mbox{练习~7.2.G}）。把
定义改成「局部看起来具有这种形状的\trRingedSpace{}之间的\trMorphism」也能行得通；
但为了避免梳理诸般细节的麻烦，我们改走另一条路——使用\trLocallyRingedSpace{}，
这相当于不仅要求函数拉回，还要求\emph{函数的值}拉回。不过我们将验证：我们最终的%
定义实际上与它等价（\mbox{练习~7.3.D}）。

我们先来正式地定义\trRingedSpace{}的\trMorphism。


\section{赋环空间的态射}

\textbf{7.2.1.\ \trhDefinition{}。}\quad \textbf{\trRingedSpace{}的\trMorphism}（morphism of ringed spaces）$\pi\colon X\to Y$ 是指：一个\trTopologicalSpace{}之间的\trContinuousMap{}（很遗憾，我们仍把它记作 $\pi$），连同一份映射 $\OO_Y\to\pi_*\OO_X$——我们把它当作「\trPullback{}映射」来想。依\trAdjointness（\S2.7.1），这与给出一个映射 $\pi^{-1}\OO_Y\to\OO_X$ 是同一回事。（把这批信息打包成 \eqref{eq:compatible-collections-diagram} 中那张图的样子，有时颇为便利。）\trMorphism{}的\trComposition{}有其显而易见的定义，故\trRingedSpace{}构成一个\trCategory。于是我们便有\trAutomorphism{}与\trIsomorphism{}的观念。你可以轻松验证：\trRingedSpace{}的\trIsomorphism{}所指与先前（定义~4.3.1）完全相同。

若 $\mathsf{U}\subset Y$ 是一个\trOpenSubset，则存在一个自然的\trRingedSpace{}态射 $(\mathsf{U},\OO_{Y|\mathsf{U}})\to(Y,\OO_Y)$（它早在 \S2.2.13 中就隐式登场过）。更精确地说：若 $\mathsf{U}\to Y$ 是 $\mathsf{U}$ 与 $Y$ 的一个\trOpenSubset{} $V$ 之间的\trIsomorphism，并且我们又给定一个\trIsomorphism{} $(\mathsf{U},\OO_{\mathsf{U}})\mathrel{\xrightarrow{\ \sim\ }}(V,\OO_{Y|V})$（经由\trIsomorphism{} $\mathsf{U}\mathrel{\xrightarrow{\ \sim\ }}V$ 而来），那么由此得到的\trRingedSpace{}映射就称为\trRingedSpace{}的一个\textbf{\trOpenEmbedding}（open embedding；或\textbf{\trOpenImmersion}，open immersion），而态射 $\mathsf{U}\to Y$ 常被写作 $\mathsf{U}\hookrightarrow Y$。

\begin{exercise}[\trRingedSpace{}的\trMorphism{}可粘]\label{ex:ringed-space-morphisms-glue}
设 $(X,\OO_X)$ 与 $(Y,\OO_Y)$ 是\trRingedSpace，$X=\cup_i U_i$ 是 $X$ 的一个\trOpenCover，而且我们有\trRingedSpace{}的\trMorphisms{} $\pi_i\colon U_i\to Y$，它们「在交叠处相合」，即 $\pi_i|_{U_i\cap U_j}=\pi_j|_{U_i\cap U_j}$。证明：存在唯一的\trRingedSpace{}态射 $\pi\colon X\to Y$ 使得 $\pi|_{U_i}=\pi_i$。（\cref{ex:morphisms-glue} 已就\trTopologicalSpace{}把这件事基本证过了。）
\end{exercise}

\begin{exercise}\label{ex:open-embeddings-monomorphisms}
证明：\trRingedSpace{}的\trOpenEmbedding{}是\trMonomorphism{}（在\trRingedSpace{}的\trCategory{}中）。
\end{exercise}

\begin{exercise}[容易的重要练习：前推 $\OO$-\trModule]\label{ex:pushforward-ox-modules}
给定\trRingedSpace{}态射 $\pi\colon X\to Y$，证明：\trSheaf{}的\trPushforward{}诱导一个\trFunctor{} $\pi_*\colon\catModOX\to\textsc{Mod}_{\OO_Y}$。
\end{exercise}

\begin{exercise}[日后要用的重要练习：拉回 $\OO$-\trModule]\label{ex:pullback-ox-modules}
（这个练习你尽可先撂在一旁，等到需要它时再去理会——因为在那之前，你未必能体会它的好处。）设 $\pi\colon(X,\OO_X)\to(Y,\OO_Y)$ 是一个\trRingedSpace{}态射。对\trInverseImage{}稍加变通，我们便能\textbf{拉回一个 $\OO_Y$-\trModule{} $\scr{G}$}，从而得到一个 $\OO_X$-\trModule{} $\pi^*\scr{G}$。
\begin{enumerate}[label=(\alph*)]
\item 证明 $(X,\pi^{-1}\OO_Y)$ 是一个\trRingedSpace，并且\trRingedSpace{}态射 $\pi$ 可分解为 $(X,\OO_X)\to(X,\pi^{-1}\OO_Y)\to(Y,\OO_Y)$。
\item 证明 $\pi^{-1}\scr{G}$ 是一个 $\pi^{-1}\OO_Y$-\trModule，并描述 $\pi^{-1}\scr{G}\otimes_{\pi^{-1}\OO_Y}\OO_X$ 如何成为一个 $\OO_X$-\trModule{}；我们称 $\pi^*\scr{G}$ 为 $\scr{G}$ 的\textbf{\trPullback}（pullback）。
\item 证明 $\pi^*$ 是一个\trCovariantFunctor{} $\textsc{Mod}_{\OO_Y}\to\catModOX$。
\item 此外，若 $\rho\colon(W,\OO_W)\to(X,\OO_X)$ 是另一个\trRingedSpace{}态射，证明「$\rho^*\pi^*=(\pi\circ\rho)^*$」（或者更精确地说：描述函子间的一个\trNaturalIsomorphism{} $\rho^*\pi^*\overset{\sim}{\longleftrightarrow}(\pi\circ\rho)^*$）。
\item（参看\cref{ex:adjointness-pushforward-inverse-image}。）证明 $\pi^*$ 是 $\pi_*$ 的\trLeftAdjoint。更精确地说：对任一 $\scr{F}\in\catModOX$，描述一个\trNatural\trBijection
\[
\Hom_{(X,\OO_X)}(\pi^*\scr{G},\scr{F})
\overset{\sim}{\longleftrightarrow}
\Hom_{(Y,\OO_Y)}(\scr{G},\pi_*\scr{F}),
\]
它在 $\scr{F}\in\catModOX$ 与 $\scr{G}\in\textsc{Mod}_{\OO_Y}$ 这两个变元上都具有\trFunctoriality。
\item（$\OO$-\trModule{}的\trPushPullMap。）设
\begin{tikzcd}[row sep=2.6em, column sep=3.6em]
W \ar[r, "\beta'"] \ar[d, "\alpha'"'] & X \ar[d, "\alpha"] \\
Y \ar[r, "\beta"] & Z
\end{tikzcd}
是\emph{\trRingedSpace}的一张\trCommutativeDiagram{}（不必为\trCartesianDiagram），而 $\scr{F}$ 是一个 $\OO_X$-\trModule。请循 \S2.7.4 对\trSheaf{}的\trPushPullMap{}的讨论，说明如何定义 $\OO_Y$-\trModule{}的\textbf{\trPushPullMap}
\[
\beta^*\alpha_*\scr{F}\longrightarrow\alpha'_*(\beta')^*\scr{F}
\]
（对照 \eqref{eq:push-pull-map}）。
\end{enumerate}
\end{exercise}

\begin{exercise}[容易的重要练习]\label{ex:stalk-map-ringed-space}
给定\trRingedSpace{}态射 $\pi\colon X\to Y$，设 $\pi(p)=q$。证明：存在\trStalk{}之间的一个映射 $(\OO_Y)_q\to(\OO_X)_p$。
\end{exercise}

\begin{exercise}[关键练习]\label{ex:ring-map-induced-morphism}
设 $\pi^{\ddagger}\colon B\to A$ 是一个\trRingMorphism。如下定义\trRingedSpace{}态射 $\pi\colon\Spec A\to\Spec B$。\trTopologicalSpace{}之间的映射已由\cref{ex:ring-maps-continuous-maps}给出。要在 $\Spec B$ 上描述\trSheaf{}态射 $\OO_{\Spec B}\to\pi_*\OO_{\Spec A}$，只需在 $\Spec B$ 的\trDistinguishedBase{}上描述一个\trSheaf{}态射。对 $D(g)\subset\Spec B$，我们定义
\[
\OO_{\Spec B}(D(g))\longrightarrow\OO_{\Spec A}(\pi^{-1}D(g))=\OO_{\Spec A}(D(\pi^{\ddagger}g))
\]
为 $B_g\to A_{\pi^{\ddagger}g}$。请验证这个定义确有意义（例如：它与 $g$ 的选取无关），并且这确实描述了\trDistinguishedBase{}上的一个\trSheaf{}态射。（这是系列练习中的第三个：\trRingMorphism{}诱导集合的映射，见 \S3.2.10；诱导\trTopologicalSpace{}的映射，见\cref{ex:ring-maps-continuous-maps}；而此刻在这里，则是\trRingedSpace{}的映射。）
\end{exercise}

关键\cref{ex:ring-map-induced-morphism}所给的\trRingedSpace{}映射其实一点也不复杂。这里就是一个例子。考虑由 $y\mapsto x^2$ 给出的\trRingMorphism{} $\mathbb{C}[y]\to\mathbb{C}[x]$（见图~\ref{fig:parabola-map}）。我们是在把带\trCoordinate{} $x$ 的\trAffineLine{}映到带\trCoordinate{} $y$ 的\trAffineLine。该映射（在\trClosedPoint{}上）是 $a\mapsto a^2$。举例来说：$[(x-3)]$ 被送到哪里？答：$[(y-9)]$，即 $3\mapsto 9$。$[(y-4)]$ 的\trInverseImage{}是什么？答：$\mathbb{C}[x]$ 中那些包含 $[(x^2-4)]$ 的\trPrimeIdeal，即 $[(x-2)]$ 与 $[(x+2)]$；所以 $4$ 的\trInverseImage{}的确是 $\pm2$。至此讲的还只是集合的映射，那已是旧闻（\S3.2.10）；现在来想想函数的拉回。$D(y-4)=\mathbb{A}^1-\{4\}$ 上的函数 $3/(y-4)$，其\trPullback{}是什么？当然是 $\mathbb{A}^1\setminus\{-2,2\}$ 上的 $3/(x^2-4)$。

关键\cref{ex:ring-map-induced-morphism}的构造很快就会化为一个\trScheme{}态射的例子。其实，我们此刻就可以把那个定义立起来。但在那之前，让我们（借下一个练习）指出：仿射\trScheme{}之间的\trRingedSpace{}态射并非个个都呈关键\cref{ex:ring-map-induced-morphism}之形态。（用 \S7.3 的语言来说：这个\trRingedSpace{}态射不是\trLocallyRingedSpace{}的态射。）%
这个练习我们并不会用到，但它回答了每个人都会有的一个疑问。

\begin{exercise}\label{ex:exotic-ringed-space-map}
回忆\cref{ex:topological-space-kx-p}：$\Spec k[y]_{(y)}$ 有两个点 $[(0)]$ 与 $[(y)]$，其中第二个是\trClosedPoint{}，第一个则不是。描述一个\emph{\trRingedSpace}的映射 $\Spec k(x)\to\Spec k[y]_{(y)}$：它把 $\Spec k(x)$ 的唯一点送到\trClosedPoint{} $[(y)]$，其在\trGlobalSections{}上的\trPullback{}映射把 $k$ 依恒等映射送到 $k$、把 $y$ 送到 $x$。证明：这个\trRingedSpace{}的映射并不呈关键\cref{ex:ring-map-induced-morphism}所述之形态。
\end{exercise}

\textbf{7.2.2.\ 暂不采用的试行定义（对照\mbox{动机~7.1.1}）。}%
\quad \trScheme{}的\trMorphism{} $\pi\colon(X,\OO_X)\to(Y,\OO_Y)$ 是这样一个\trRingedSpace{}态射：它「在局部看似」关键\cref{ex:ring-map-induced-morphism}所描述的那种仿射\trScheme{}映射。确切地说：对每一组\trAffine\trOpenSet{} $\Spec A\subset X$、$\Spec B\subset Y$，只要 $\pi(\Spec A)\subset\Spec B$，其所诱导的\trRingedSpace{}映射就应当呈关键\cref{ex:ring-map-induced-morphism}所示之形态。

我们希望这个定义在任一\trAffine\trCover{}上都可检验，而且我们本可指望借\trAffineCommunicationLemma{}按这条路子把理论发展起来。这也行得通，不过更便利的是来一记更巧的妙招：在下一节中，我们会用一用\trLocallyRingedSpace{}这个观念；而一旦用过，就会把它像昨天的垃圾一样扔掉。%


\section{从\trLocallyRingedSpace{}到概形的态射}

为了证明\trMorphism{}确会如我们所望的那样行事，我们将动用\trLocallyRingedSpace{}这一观念。虽然它在几何学的其他场合自有用处，但往后我们不会再用到它。\trLocallyRingedSpace（连同其间的映射）这一观念，其灵感来自我们对\trManifold{}的已有认识（见\cref{ex:stalk-map-preserves-vanishing}）。设 $\pi\colon X\to Y$ 是\trManifold{}之间的一个\trMorphism，$\pi(p)=q$，而 $f$ 是 $Y$ 上一个在 $q$ 处为零的函数，那么拉回到 $X$ 上的函数 $\pi^{\sharp}(f)$ 就应当在 $p$ 处为零。换个说法：在 $q$ 处为零的函数\trGerm（于 $q\in Y$ 处取的），拉回后应当成为在 $p$ 处为零的函数\trGerm（于 $p\in X$ 处取的）。

\textbf{7.3.1.}\ \emph{\trhDefinition{}。}\quad 回想（定义~4.3.7）：\trLocallyRingedSpace{}是这样的\trRingedSpace{} $(X,\OO_X)$，其诸\trStalk{} $\OO_{X,p}$ 全都是\trLocalRing。\textbf{\trLocallyRingedSpace{}之间的\trMorphism} $\pi\colon X\to Y$，是这样一个\trRingedSpace{}的\trMorphism：它所诱导的诸\trStalk{}之间的映射 $\pi^{\sharp}\colon\OO_{Y,q}\to\OO_{X,p}$（练习~7.2.E）把 $\OO_{Y,q}$ 的\trMaximalIdeal{} $\mathfrak{m}_{Y,q}$ 送入 $\OO_{X,p}$ 的\trMaximalIdeal{} $\mathfrak{m}_{X,p}$（即一个「\trLocalRing{}的\trRingMorphism」（morphism of local rings））。（特别提请注意：\trLocallyRingedSpace{}构成一个\trCategory。）这话说的是一件相当具体而直观的事情：「若 $p\mapsto q$，而 $g$ 是一个在 $q$ 处为零的函数，那么它拉回后便是一个在 $p$ 处为零的函数。」你还会想要：「若 $p\mapsto q$，而 $g$ 是一个在 $q$ 处\emph{不}为零的函数，那么它拉回后便是一个在 $p$ 处\emph{不}为零的函数。」而这正是下一道练习的一个推论。

\textbf{7.3.A.（练习：\trLocallyRingedSpace{}上函数的值可拉回）}\quad 设 $\pi\colon X\to Y$ 是\trLocallyRingedSpace{}之间的一个\trMorphism，且 $\pi(p)=q$。证明：$\pi$ 诱导诸\trResidueField{}之间的一个包含映射 $\kappa(q)\hookrightarrow\kappa(p)$。这就抓住了如下事实：拉回把「函数为零（相应地，非零）的\trLocus」送到「拉回后的函数为零（相应地，非零）的\trLocus」。特别地，由\cref{ex:vanishing-locus-closed}(b)，可逆函数的拉回仍是可逆函数。

小结一下：\trLocallyRingedSpace{}这个观念，我们只在定义\trScheme{}的\trMorphism{}并证明\trMorphism{}具有种种合理性质时才会用到。关于\trLocallyRingedSpace，你需要记住的要点有二：(i) 函数在点处有值；(ii) 给定\trLocallyRingedSpace{}之间的一个映射，函数的值会「拉回」，而且特别地，函数的零点也会「拉回」。

\textbf{7.3.B.（练习）}\quad 证明：\trLocallyRingedSpace{}的\trMorphism{}可粘（对照练习~7.2.A）。（提示：你对练习~7.2.A 的解答或许原封不动就能用。）

\textbf{7.3.C.（容易的重要练习）}\quad 回想（\S4.3.7）：$\Spec A$ 是一个\trLocallyRingedSpace。证明：由\trRingMorphism{} $\pi^{\sharp}\colon B\to A$ 所定义的\trRingedSpace{}的\trMorphism{} $\pi\colon\Spec A\to\Spec B$（练习~7.2.F）是\trLocallyRingedSpace{}之间的一个\trMorphism。

\textbf{7.3.2.}\ \emph{\textbf{关键命题。}}\ ——\ \emph{若 $\pi\colon\Spec A\to\Spec B$ 是\trLocallyRingedSpace{}之间的一个\trMorphism，那么它正是由映射 $\pi^{\sharp}\colon B=\Gamma(\Spec B,\OO_{\Spec B})\to\Gamma(\Spec A,\OO_{\Spec A})=A$ 依练习~7.3.C 之例所诱导的那个\trLocallyRingedSpace{}之间的\trMorphism。}

（旁注：\cref{ex:scheme-isos-ring-isos} 是关键命题~7.3.2 的一个特例。你应当回头看看自己对\cref{ex:scheme-isos-ring-isos}的解答，找一找自己在何处暗中用到了关于\trLocallyRingedSpace{}的想法。）

\emph{证明。}\ 设 $\pi\colon\Spec A\to\Spec B$ 是\trLocallyRingedSpace{}之间的一个\trMorphism。我们要证明：它被其在\trGlobalSections{}上的映射 $\pi^{\sharp}\colon B\to A$ 完全确定。首先需要核验：点的映射由\trGlobalSections{}确定。$\Spec A$ 的一个点 $p$，可以等同于「在其上为零的整体函数所成的那个\trPrimeIdeal」。依练习~7.3.A，作为一个\trIdeal{}的 $\pi(p)$ 就是自然映射 $B\to\kappa(p)$ 的\trKernel。（这里我们用到了 $\pi$ 是\trLocallyRingedSpace{}之间的映射这一事实。）这正是「由环映射 $B\to A$ 诱导的集合映射 $\Spec A\to\Spec B$」当初被定义的方式（\S3.2.10）。

特别提请注意：若 $g\in B$，则 $\pi^{-1}(D(g))=D(\pi^{\sharp}g)$——这又一次用到了 $\pi$ 是\trLocallyRingedSpace{}的\trMorphism{}这一假设。

尚待证明的是：$\pi^{\sharp}\colon\OO_{\Spec B}\to\pi_*\OO_{\Spec A}$ 正是练习~7.2.F 所给出的那个\trSheaf{}的\trMorphism（对照练习~7.3.C）。在特异基底上核验它就足够了（\cref{ex:sheaf-morphisms-on-base}(a)）。现在我们想核验：对任何一个诱导出\trSheaf{}的映射 $\OO_{\Spec B}\to\pi_*\OO_{\Spec A}$ 的\trLocallyRingedSpace{}之间的映射，它在任一\trDistinguishedOpen{} $D(g)\subset\Spec B$（$g\in B$）上的截面映射，都由\trGlobalSections{}的映射 $B\to A$ 所确定。

考虑交换图
\[
\begin{tikzcd}[row sep=3.4em]
B \ar[r, equal] \ar[d, "\mathrm{res}_{\Spec B,\,D(g)}"] &
\Gamma(\Spec B,\OO_{\Spec B}) \ar[r, "\pi^{\sharp}{}_{\Spec B}"] &
\Gamma(\Spec A,\OO_{\Spec A}) \ar[r, equal] \ar[d, "\mathrm{res}_{\Spec A,\,D(\pi^{\sharp}g)}"'] &
A \\
B_g \ar[r, equal] &
\Gamma(D(g),\OO_{\Spec B}) \ar[r, "\pi^{\sharp}{}_{D(g)}"] &
\Gamma(D(\pi^{\sharp}g),\OO_{\Spec A}) \ar[r, equal] &
A_{\pi^{\sharp}g}=A\otimes_B B_g
\end{tikzcd}
\]
竖直的箭头（向诸\trDistinguishedOpen{}的限制）都是对 $g$ 的\trLocalization，于是下方的水平映射 $\pi^{\sharp}{}_{D(g)}$ 便由上方的映射所确定（它无非是对 $g$ 作\trLocalization）。\hfill$\square$

我们已经为下面的定义做好了准备。

\textbf{7.3.3.}\ \emph{\trhDefinition{}。}\quad 若 $X$ 与 $Y$ 是\trScheme，那么作为\trLocallyRingedSpace{}的\trMorphism{} $\pi\colon X\to Y$ 就称为一个\textbf{\trScheme{}的\trMorphism}（morphism of schemes）。至此我们便定义了\textbf{\trScheme{}的\trCategory}，记作 $\mathit{Sch}$。（于是我们也就有了\trIsomorphism{}的观念——与先前 \S4.3.1 处的含义相同——以及\trAutomorphism{}的观念。当我们把 $\pi$ 解读为「由 $Y$ 参数化的一族概形」时，$\pi$ 的\emph{\trTarget} $Y$ 有时被称为\textbf{\trBase\trScheme}（base scheme）或\textbf{\trBase}（base）；待我们在 \S10.3.2 中定义了\trMorphism{}的诸\trFiber{}之后，这话会更加清楚。）

用\trLocallyRingedSpace{}陈述的这条定义，很容易推出试行定义~7.2.2：

\textbf{7.3.D.（重要练习）}\quad 证明：\trScheme{}的\trMorphism{} $\pi\colon X\to Y$ 是一个「在局部看似诸\trAffineScheme{}的\trMorphism」的\trRingedSpace{}的\trMorphism。确切地说：若 $\Spec A$ 是 $X$ 的一个\trAffine\trOpenSubset，$\Spec B$ 是 $Y$ 的一个\trAffine\trOpenSubset，且 $\pi(\Spec A)\subset\Spec B$，则其所诱导的\trRingedSpace{}的\trMorphism{}是一个\trAffineScheme{}的\trMorphism。（为助你一臂之力，请注意：若 $W\subset X$ 与 $Y\subset Z$ 都是\trRingedSpace{}的\trOpenEmbedding，则任何\trRingedSpace{}的\trMorphism{} $X\to Y$ 都经复合 $W\to X\to Y\to Z$ 诱导出\trRingedSpace{}的一个\trMorphism{} $W\to Z$。）证明：只需在一组 $(\Spec A_i,\Spec B_i)$ 上核验即可，这里诸 $\Spec A_i$ 构成 $X$ 的一个\trOpenCover，诸 $\Spec B_i$ 构成 $Y$ 的一个\trOpenCover。

在实际操作中，我们将使用仿射覆盖式的解读，而把\trLocallyRingedSpace{}彻底抛诸脑后。特别地，说不严格一点：\trAffineScheme{}的\trCategory{}就是把箭头全部掉头的\trRing{}的\trCategory。更精确地说：

\textbf{7.3.E.（练习）}\quad 证明：\trRing{}的\trCategory{}与\trAffineScheme{}的\trOppositeCategory{}是等价的（关于\trOppositeCategory{}与\trEquivalenceOfCategories，可分别参看 \S1.1.16 与 \S1.1.21）。

特别地，从一个点到一个点的映射可以有许多个！比如，从点 $\Spec\mathbb{C}$ 到点 $\Spec\mathbb{C}$ 就有两个不同的映射：恒等映射（对应于恒等映射 $\mathbb{C}\to\mathbb{C}$），以及复共轭。（这样的映射甚至还有更多！）

显然（由关于\trLocallyRingedSpace{}的相应事实可知）：\trScheme{}的\trMorphism{}可粘（练习~7.3.B），且两个\trMorphism{}的\trComposition{}仍是\trMorphism。这个\trCategory{}中的\trIsomorphism{}，恰恰就是我们先前（\S4.3.1）定义它们时的样子。

\textbf{7.3.F.（富于启发的练习）}\quad （这道练习可以让你练一练「把概形切成\trAffine\trOpenSet{}来理解\trMorphism」的手艺。）请把下面这句话琢磨出道理来：「$\mathbb{A}^{n+1}_k\setminus\{\tilde{0}\}\to\mathbb{P}^n_k$ 由
\[
(x_0,x_1,\dots,x_n)\longmapsto[x_0,x_1,\dots,x_n]
\]
给出，是一个\trScheme{}的\trMorphism。」警示：你不能只说点往哪里去，还得说函数往哪里去。所以你也许得把 $\mathbb{A}^{n+1}_k\setminus\{\tilde{0}\}$ 切成一块块仿射的，逐一描述那些映射，并核验它们可粘。（把 $k$ 换成一般的\trRing{} $B$ 时，你还能推广吗？答案见\cref{ex:no-common-zeros-projective}。）

\textbf{7.3.4.\ 到\trAffineScheme{}的\trMorphism。}\quad 下面这项结果表明：要描述到一个\trAffineScheme{}的\trMorphism，无须费力气用诸\trAffine\trOpenSet{}把源头盖满，这很容易。

\textbf{7.3.G.（必作练习）}\quad 证明：诸\trMorphism{} $X\to\Spec A$ 与诸\trRingMorphism{} $A\to\Gamma(X,\OO_X)$ 天然地一一对应。提示：先证 $X$ 为\trAffine{}时如此；再用\trMorphism{}可粘这一事实（练习~7.3.B）。（这话在\trLocallyRingedSpace{}的\trCategory{}中甚至也成立。你尽可以在那样的一般性下证明它，但在\trScheme{}的\trCategory{}中要更容易。你或许会在这句更一般的陈述里，听到双星注记~6.2.5 那句格言的回声。）

特别地，从每个\trScheme{}到「其\trGlobalSections{}之\trRing{}的 $\Spec$」都存在一个\trCanonical{}的\trMorphism。（警告：即便 $X$ 是有限型的 $k$-\trScheme，\trGlobalSections{}之\trRing{}也可能很惹人厌！特别地，它未必\trFinitelyGenerated；见 \S19.11.13。）\trCanonical{}态射 $X\to\Spec\Gamma(X,\OO_X)$ 是\trIsomorphism，当且仅当 $X$ 是\trAffine{}的（即：\trIsomorphic{}于某个 $\Spec A$，其中 $A$ 是\trRing）；而且在此时，它正是注记~4.3.2 所暗示的那个\trIsomorphism。

\textbf{7.3.H.（容易的练习）}\quad 若 $S_\bullet$ 是一个\trFinitelyGenerated{}的分次 $A$-\trAlgebra，请描述一个自然的「\trStructureMorphism{}」（structure morphism）$\Proj S_\bullet\to\Spec A$。%

\textbf{7.3.5.}\ \emph{\trhRemark{}。}\quad 从必作练习~7.3.G 出发，只消一小步便能证明：概形的某些\trProduct{}是存在的：若 $A$ 与 $B$ 是\trRing，则 $\Spec A\times\Spec B=\Spec(A\otimes_{\mathbb{Z}}B)$；而若 $A$ 与 $B$ 是 $\mathbb{C}$-\trAlgebra，则 $\Spec A\times_{\Spec\mathbb{C}}\Spec B=\Spec(A\otimes_{\mathbb{C}}B)$。不过我们并不着急，且等到\mbox{练习~10.1.A}再正经讨论此事。

\textbf{7.3.6.}\ \extrahard\ \emph{供行家的旁注：$\Gamma$ 与 $\Spec$ 是\trAdjoint{}的。}\quad 我们有一个从\trRing{}到\trLocallyRingedSpace{}的\trContravariantFunctor{} $\Spec$，以及一个从\trLocallyRingedSpace{}到\trRing{}的\trContravariantFunctor{} $\Gamma$。事实上，$(\Gamma,\Spec)$ 竟是一个\trAdjointPair！（警示：我们只讨论过\trCovariantFunctor{}的伴随；倘若你当真在意，就得自己琢磨出如何为\trContravariantFunctor{}定义伴随。）如此一来，我们本可以通过径直要求 $\Spec$ 是 $\Gamma$ 的\trRightAdjoint{}来定义 $\Spec$。（一个有趣却不相干的旁问：如果你用的是\trRingedSpace{}而非\trLocallyRingedSpace，$\Gamma$ 依然有一个\trRightAdjoint。它是什么？）

我们能够轻易描述驶向\trAffineScheme{}的\trMorphism，这将使我们得以重访先前关于给定域或环上的概形的讨论（\S5.3.6）。

\textbf{7.3.7.\ \trComplex{}概形的范畴（或更一般地，$k$ 为\trField{}时诸 $k$-概形的范畴；或更一般地，$A$ 为\trRing{}时诸 $A$-概形的范畴；或更一般地，$S$ 为\trScheme{}时诸 $S$-概形的范畴）。}\quad $S$-\trScheme{}的\trCategory{} $\mathit{Sch}_S$（其中 $S$ 是一个\trScheme）定义如下。其\trObject（诸 $S$-概形）是形如
\[
\begin{tikzcd}[row sep=2.4em]
X \ar[d] \\
S
\end{tikzcd}
\]
的\trMorphism。（到 $S$ 的那个\trMorphism{}称为\textbf{\trStructureMorphism{}}。这一叫法的一个动机是：若 $S=\Spec A$，则\trStructureMorphism{}让 $X$ 的每个\trOpenSet{}上的函数都带有 $A$-\trAlgebra{}的结构；参看 \S5.3.6。）$S$-概形的\trCategory{}中的诸\trMorphism{}定义为交换图
\[
\begin{tikzcd}[row sep=2.4em, column sep=4em]
X \ar[r] \ar[d] & Y \ar[d] \\
S \ar[r, "="] & S
\end{tikzcd}
\]
它们更便利的写法是
\[
\begin{tikzcd}[row sep=2.4em, column sep=4em]
X \ar[r] \ar[dr] & Y \ar[d] \\
& S
\end{tikzcd}
\]
只要不引起混淆（即\trBase\trScheme{}已经自明），就只写出图中上面那一行。当 $S=\Spec A$（$A$ 为\trRing）时，我们便得到 $A$-\trScheme{}的观念；它与 \S5.3.6 中的那个定义是同一回事（练习~7.3.I），只是形式更令人满意。例如，复几何学家便可以考虑 $\mathbb{C}$-概形的范畴。

\textbf{7.3.I.（容易的练习）}\quad 证明：\S7.3.7 中给出的这个 $A$-\trScheme{}的定义，与 \S5.3.6 中较早的定义相一致。

\textbf{7.3.J.（容易的练习）}\quad 证明：$\Spec\mathbb{Z}$ 是\trScheme{}的\trCategory{}中的\trFinalObject。换言之，若 $X$ 是任何一个\trScheme，则存在唯一的到 $\Spec\mathbb{Z}$ 的\trMorphism。（因此，\trScheme{}的\trCategory{}同构于 $\mathbb{Z}$-\trScheme{}的\trCategory。）若 $A$ 是任何一个\trRing（例如某个\trField{} $k$），证明：$\Spec A$ 是 $A$-\trScheme{}的\trCategory{}中的\trFinalObject。

\textbf{7.3.8.\ \trScheme{}的\trCoproduct。}

\textbf{7.3.K.（容易的练习）}\quad 设 $X$ 与 $Y$ 是\trScheme。证明：二者的\trDisjointUnion{} $X\sqcup Y$（定义见\cref{ex:disjoint-union-schemes}）是\trScheme{}的\trCategory{}中 $X$ 与 $Y$ 的\trCoproduct，从而 \trCoproduct{}符号 $\sqcup$ 的使用便有了着落。（与\trIdempotent{}的一个出人意料的关联，见 \S3.6.3。）

\textbf{7.3.9.\ 从（某些）\trAffineScheme{}出发的\trMorphism。}\quad \emph{从}\trAffineScheme{}出发的\trMorphism，不像\emph{到}\trAffineScheme{}的\trMorphism{}那么简单，但有几种情形值得一提。

\vkexsync{L}
\begin{exercise}\label{ex:spec-local-ring-canonical-map}
(a)~设 $p$ 是\trScheme{} $X$ 的一个点。请描述一个\trCanonical{}（即无须选取的）\trMorphism{} $\Spec\OO_{X,p}\to X$。（提示：先对\trAffine{}的 $X$ 做这件事。但随后对一般的 $X$，务必证明你的\trMorphism{}与选取无关。）
(b)~定义一个\trCanonical{}的\trMorphism{} $\Spec\kappa(p)\to X$。（这个态射常被写作 $p\to X$；此时我们给 $p$ 一种显而易见的解读：把它看作\trScheme{} $\Spec\kappa(p)$。）
\end{exercise}

\textbf{7.3.M.（练习：从局部环的谱到 $X$ 的\trMorphism）}\quad 设 $X$ 是一个\trScheme，$(A,\mathfrak{m})$ 是一个\trLocalRing。设我们有一个\trScheme{}的\trMorphism{} $\pi\colon\Spec A\to X$，它把 $[\mathfrak{m}]$ 送到 $p$。证明：任何包含 $p$ 的\trOpenSet{}都包含 $\pi$ 的像。证明：$\Mor(\Spec A,X)$ 与
\[
\bigl\{(p\in X,\ \text{某个局部\trHomomorphism{}}\ \OO_{X,p}\to A)\bigr\}
\]
之间成双射。（可能的提示：\cref{ex:spec-local-ring-canonical-map}(a)。）

这些练习把我们引向\textbf{域\trValuedPoint{}}（field-valued points）的观念，进而是一般些的\textbf{环\trValuedPoint{}}（ring-valued points）的观念，再进而是一般些的\textbf{概形\trValuedPoint{}}（scheme-valued points）的观念。%

\textbf{7.3.10.\ \trhDefinition{}：\trFunctorOfPoints，与\trScheme{}的概形\trValuedPoint{}（以及环\trValuedPoint{}与域\trValuedPoint{}）。}\quad 若 $Z$ 是一个\trScheme，则\trScheme{} $X$ 的\textbf{$Z$-\trValuedPoint{}}（$Z$-valued points；记作 $X(Z)$）定义为诸映射 $Z\to X$。若 $A$ 是一个\trRing，则\trScheme{} $X$ 的\textbf{$A$-\trValuedPoint{}}（$A$-valued points；记作 $X(A)$）定义为该概形的诸 $(\Spec A)$-\trValuedPoint{}。（最常见的情形是 $A$ 为一个\trField。）

如果你在某个基底概形 $B$ 之上工作——例如，复代数几何学家将只考虑 $B=\Spec\mathbb{C}$ 上的概形与态射——那么在上面的定义里就暗含了一张结构映射 $Z\to B$（在 $X(A)$ 的情形则是 $\Spec A\to B$）。例如，对一位复几何学家而言，若 $X$ 是 $\mathbb{C}$ 上的概形，则 $X$ 的 $\mathbb{C}(t)$-\trValuedPoint{}对应于形如
\[
\begin{tikzcd}[row sep=2.6em, column sep=4em]
\Spec\mathbb{C}(t) \ar[r] \ar[dr, "\xi"'] & X \ar[d, "\pi"] \\
& \Spec\mathbb{C}
\end{tikzcd}
\]
的交换图，其中 $\pi\colon X\to\Spec\mathbb{C}$ 是 $X$ 的结构映射，而 $\xi$ 对应于环的显然包含 $\mathbb{C}\to\mathbb{C}(t)$。（警告：$k$-\trScheme{} $X$ 的一个 $k$-\trValuedPoint{}有时被称为 $X$ 的一个「\textbf{\trRationalPoint{}}」（rational point），这个叫法很危险，因为对世间大多数人而言，「有理」指的是 $\mathbb{Q}$。我们将使用更稳妥的说法：$X$ 的「$k$-\trValuedPoint{}」。另一个稳妥的选择是 $X$ 的「$k$-\trRationalPoint{}」或 $k$-点。）

「$Z$-\trValuedPoint」（以及 $A$-\trValuedPoint）这一叫法并不妥帖，因为我们先前已经定义过\trScheme{}的「点」这个观念，而 $Z$-\trValuedPoint（以及 $A$-\trValuedPoint）却（未必）是点！但这类用法在文献中早已根深蒂固。（不妨到索引的「点」（point）条目下去看看，修饰这个词的形容词还有更多彼此打架的用法。）

\textbf{7.3.11.}\quad 此外，我们将会看到「$Z$-\trValuedPoint{}的\trProduct」会如你所望的那样行事（\S10.1.3）。这套语言之所以能发人联想，还有一层相关的原因：记号 $X(Z)$ 暗示着把 $X$ 解读为一个从\trScheme{}到\trSet{}的\textbf{\trContravariantFunctor} $h_X$——即\trScheme{} $X$ 的（概形值）\trFunctorOfPoints（参看例~1.1.20）。

$A$-\trValuedPoint{}之所以有用，还有一个下里巴人式的理由：\emph{对诸 $f_i\in\mathbb{Z}[x_1,\dots,x_n]$（关系）而言，仿射概形 $\Spec\mathbb{Z}[x_1,\dots,x_n]/(f_1,\dots,f_r)$ 的 $A$-\trValuedPoint{}恰是方程组}
\[
f_1(x_1,\dots,x_n)=\cdots=f_r(x_1,\dots,x_n)=0
\]
\emph{在\trRing{} $A$ 中的解。}例如，$x^2+y^2=16$ 的有理解，恰是 $\Spec\mathbb{Z}[x,y]/(x^2+y^2-16)$ 的 $\mathbb{Q}$-\trValuedPoint{}；整数解恰是 $\mathbb{Z}$-\trValuedPoint{}。于是，仿射概形（在 $\mathbb{Z}$ 上\trFiniteType{}的）的 $A$-\trValuedPoint{}可以得到简单的解读。而在 $A$ 为局部环的特殊情形，一般概形的 $A$-\trValuedPoint{}也有很好的解读（\cref{ex:spec-local-ring-canonical-map} 与练习~7.3.M）。

\vkexsync{N}
\begin{exercise}\label{ex:valued-points-determine-morphisms}
(a)~（容易）证明：\trScheme{}的\trMorphism{} $X\to Y$ 诱导诸 $Z$-\trValuedPoint{}之间的一个映射 $X(Z)\to Y(Z)$。
(b)~注意：\trScheme{}的\trMorphism{} $X\to Y$ 并不由其「底下的」点的映射确定。（例子何在？）证明：它们\emph{确实}由其诱导的诸 $Z$-\trValuedPoint{}的映射所确定，其中 $Z$ 跑遍一切\trScheme。（提示：取 $Z=X$。在做这道练习的过程中，你将借着\mbox{练习~10.1.B}的外衣，把\trYonedaLemma{}大体证明出来。）
\end{exercise}

另一方面，\trProjectiveSpace{}的 $Z$-\trValuedPoint{}可能颇为微妙。有些映射我们能轻易写出——把下一道练习应用于 $X=\Spec A$（其中 $A$ 是一个 $B$-\trAlgebra）的情形，便可见一斑。

\begin{exercise}[容易（但出奇富于启发）的练习（参看练习~7.3.F）]\label{ex:no-common-zeros-projective}
(a)~设 $B$ 是一个\trRing。若 $X$ 是一个 $B$-\trScheme，而 $f_0,\dots,f_n$ 是 $X$ 上没有公共零点的 $n+1$ 个函数，证明：$[f_0,\dots,f_n]$ 给出一个 $B$-\trScheme{}的\trMorphism{} $X\to\mathbb{P}^n_B$。
(b)~设 $g$ 是 $X$ 上一个无处为零的函数，$f_i$ 如 (a) 中所述。证明：到 $\mathbb{P}^n_B$ 的诸\trMorphism{} $[f_0,\dots,f_n]$ 与 $[gf_0,\dots,gf_n]$ 是同一个。
\end{exercise}

\textbf{7.3.12.}\ \emph{\trhExample{}（重访练习~7.3.F）。}\quad 考虑 $\mathbb{A}^{n+1}$ 上的 $n+1$ 个函数 $x_0,\dots,x_n$（又名\trTrivialBundle{}（trivial bundle）的 $n+1$ 个\trSection）。它们在 $\mathbb{A}^{n+1}\setminus\{\tilde{0}\}$ 上没有公共零点。于是它们确定一个\trMorphism{} $\mathbb{A}^{n+1}\setminus\{\tilde{0}\}\to\mathbb{P}^n$。%

\textbf{7.3.13.}\quad 你或许会指望\cref{ex:no-common-zeros-projective}(a) 给出到射影空间（$B$ 上的）的一切\trMorphism。但事实并非如此。的确，就连恒等\trMorphism{} $X=\mathbb{P}^1_k\to\mathbb{P}^1_k$ 也不是这种形状，因为源头 $\mathbb{P}^1$ 手里没有非常值的\trGlobalSections{}可用来搭出这个映射。（源头为仿射的类似例子也有。）不过，\mbox{定理~15.2.2} 中有一个正确的推广（它刻画了从概形到射影概形的\emph{一切}映射）。这条结果粗略地讲：只要诸 $f_i$ 不是地道的函数，而是某个\trLineBundle{}的\trSection，这套想法就行得通。想要干干净净地理解到射影概形的映射，将成为理解\trLineBundle{}的一路重要动机。

下一节中我们将会见到更多描述到射影空间的映射的途径。\mbox{练习~15.2.A} 将给出另一种描述，它直接推广了\cref{ex:no-common-zeros-projective}(a)，并且将（在\mbox{定理~15.2.2} 中）被证明是一种「泛」的描述。

顺带一提：在\trGrothendieck{}之前，「该如何正确解读点带有环中『坐标』的射影空间」曾被当作一个实实在在的问题。这些困难源于函子式推理的缺席。而通往正确答案的线索其实早已存在（从\trManifold{}到 $\mathbb{RP}^n$ 的映射也会遭遇同样的问题）——如果你提出的是一个这样的几何问题（射影空间可是几何对象），答案就必然是几何的，而非纯代数的！

\textbf{7.3.14.\ \trMorphism{}的可视化：当\trNilpotent{}在场时为概形的映射绘图。}\quad 你现在已经知道如何让\trScheme{}的诸点（\S3.3）与\trNilpotent（\S4.2 与 \S6.6）在眼前显形。下面这道不精确的练习，将让你对「当\trNilpotent{}掺和进来时，如何为概形的映射绘图」有一点感觉。设 $a\in\mathbb{C}$。考虑由 $x\mapsto a\epsilon$ 给出的环映射 $\mathbb{C}[x]\to\mathbb{C}[\epsilon]/(\epsilon^2)$。回想：$\Spec\mathbb{C}[\epsilon]/(\epsilon^2)$ 可以画作一个带着\trTangentVector{}（tangent vector）的点（\S4.2）。若 $a\neq0$，你会如何为这个映射画像？若 $a=0$，你的图景又会如何改变？（此时\trTangentVector{}应当被「碾碎」。）

\mbox{练习~13.1.G} 将把这一点大幅推广。你不妨现在就去读读它的陈述，以期一乐。

\textbf{7.3.15.}\ \extrahard\ \textbf{供有相应背景的读者：\trComplex{}代数簇的\trAnalytification{}（analytification）。}\quad 警告：任何关于\trAnalytification{}的讨论，都只面向熟悉复解析簇这一观念的读者，以及有意愿沿着我们对概形的展开、自行平行地把它发展出来的读者。%

\textbf{7.3.P.（练习：把\trComplex{}代数簇\trAnalytification{}）}\quad 设 $X$ 是一个既约的有限型 $\mathbb{C}$-\trScheme。请定义相应的\textbf{复解析\trPrevariety{}}（complex analytic prevariety）$X_{\mathrm{an}}$。（解析\trPrevariety{}的定义，就是去掉\trHausdorff{}条件的\trVariety{}的定义。）警示：你的定义不得依赖于对 $X$ 的某个\trAffine\trCover{}的选取。（提示：先讲清楚如何把既约的有限型仿射 $\mathbb{C}$-概形\trAnalytification{}；然后粘合。）利用弱\trNullstellensatz（\cref{thm:hilbert-weak-nullstellensatz}），给出 $X$ 的诸\trClosedPoint{}与 $X_{\mathrm{an}}$ 的诸点之间的一个双射。（事实上，还可以构造出诸集合之间的一个\trContinuousMap{} $X_{\mathrm{an}}\to X$，它是\cref{ex:map-c2-a2q} 的推广。）%

\textbf{7.3.Q.（练习：\trAnalytification{}函子）}\quad 回想练习~7.3.P 中的\trAnalytification{}构造。对既约有限型 $\mathbb{C}$-概形的每个\trMorphism{} $\pi\colon X\to Y$（$\mathbb{C}$ 上的），定义复解析\trPrevariety{}之间的一个\trMorphism{} $\pi_{\mathrm{an}}\colon X_{\mathrm{an}}\to Y_{\mathrm{an}}$（即 $\pi$ 的\trAnalytification{}）。证明：\trAnalytification{}给出一个从既约有限型 $\mathbb{C}$-概形的\trCategory{}到复解析\trPrevariety{}的\trCategory{}的\trFunctor。

\textbf{7.3.16.}\ \emph{\trhRemark{}。}\quad 两个互不同构的\trVariety{}可以拥有同构的\trAnalytification{}。例如，\trSerre{}在复\trManifold{} $\mathbb{C}^*\times\mathbb{C}^*$ 上描述过两种不同的代数结构（见 [Ha2, p.~232] 与 [MO68421]\footnote{译注：[MO68421] 指 MathOverflow 问答第 68421 号。}）：一种是「显然的那一种」，另一种则是某条\trEllipticCurve{}上的一个 $\mathbb{P}^1$-丛再去掉一个\trSection。至于带有无穷多个代数结构的光滑复曲面，可参看 \S19.11.3。另一方面，紧致的复代数簇只能有一种代数结构（见 [Se3, \S19]）。

\textbf{7.3.17.}\ \emph{关于\trAnalytification{}的进一步事实。}\quad 欲知\trAnalytification{}的更多内容，可参看练习~11.2.K（\trHausdorff{}性）、\mbox{练习~14.1.J}、\mbox{练习~18.4.H}、\mbox{练习~18.4.U}（线丛）、\S18.6.5、\S20.1.9（次数）与\mbox{练习~21.2.W}（\trManifold）。复解析空间的背景，见 [GR]。


\section{分次环的映射与射影概形的映射}

正如\trRing{}的\trMap{}以相反的方向对应于\trAffineScheme{}的\trMap，（\trBaseRing{} $A$ 上的）\trGradedRing{}的\trMap{}有时也以相反的方向给出\trProjectiveScheme{}的\trMap。这是一次并不完美的推广：并非每个\trGradedRing{}的\trMap{}都给出\trProjectiveScheme{}的\trMap（\S7.4.2）；并非每个\trProjectiveScheme{}的\trMap{}都来自\trGradedRing{}的\trMap（\S19.11.9）；而且不同的\trGradedRing{}的\trMap{}可以给出同一个\trScheme{}的\trMap（练习~7.4.C）。

在做下面这道练习时，把例~7.4.1 与例~7.4.2 想上一遍，或许不无助益。

\textbf{7.4.A.（必作练习）}\quad 设 $\phi\colon S_\bullet\longrightarrow R_\bullet$ 是 ($\mathbb{Z}^{\ge0}$-)\trGradedRing{}之间的一个\trMorphism。（所谓\textbf{\trGradedRing{}映射}（map of graded rings），我们指的是这样一个\trRing{}的\trMap：它作为「\trGrading{}\trAbelianSemigroup」（或「\trGrading{}\trCommutativeMonoid」）之间的\trMap{}保持\trGrading。换言之，存在某个 $d>0$，使得对一切 $n$，$S_n$ 都被映入 $R_{dn}$。）证明：它诱导\trScheme{}的一个\trMorphism
\[
(\Proj R_\bullet)\setminus V(\phi(S_+))\longrightarrow\Proj S_\bullet.
\]
（提示：设 $f$ 是 $S_+$ 的一个\trHomogeneous{}元素。定义一个\trMap{} $D_+(\phi(f))\to D_+(f)$。证明它们可以粘合起来（当 $f$ 跑遍 $S_+$ 的全部\trHomogeneous{}元素时）。证明这就定义了一个从整个 $\Proj R_\bullet\setminus V(\phi(S_+))$ 出发的\trMap。）特别地，若
\begin{equation*}\tag{7.4.0.1}\label{eq:vphi-s-plus-empty}
V(\phi(S_+))=\varnothing,
\end{equation*}
则我们便有\trMorphism{} $\Proj R_\bullet\to\Proj S_\bullet$。从你的解答可以看出：若 $\phi$ 更是 $A$-\trAlgebra{}之间的\trMorphism，那么所诱导的\trMorphism{} $\Proj R_\bullet\setminus V(\phi(S_+))\to\Proj S_\bullet$ 就是 $A$-\trScheme{}之间的\trMorphism。

\textbf{7.4.1.}\ \emph{例。}\quad 让我们看看练习~7.4.A 的实际运用。我们将以概形论的方式解读复射影流形之间的如下\trMap
\[
\mathbb{CP}^1\longrightarrow\mathbb{CP}^2
\]
\[
[s,t]\longmapsto[s^{20},s^9t^{11},t^{20}].
\]
首先注意，这是\trWellDefined{}的：$[\lambda s,\lambda t]$ 与 $[s,t]$ 被送到 $\mathbb{CP}^2$ 中同一个点。它之所以\trWellDefined，缘由在于：$s^{20}$、$s^9t^{11}$、$t^{20}$ 这三个多项式都是次数 $20$ 的\trHomogeneous{}多项式。

在代数上，这对应于一个方向相反的\trGradedRing{}的\trMap
\[
\mathbb{C}[x,y,z]\longrightarrow\mathbb{C}[s,t],
\]
由 $x\mapsto s^{20}$, $y\mapsto s^9t^{11}$, $z\mapsto t^{20}$ 给出。你应当依自己对练习~7.4.A 的解答来解读这件事，并与 \S3.2.11 的仿射例子相比较。

\textbf{7.4.2.}\ \emph{例。}\quad 注意，并不存在由 $[x,y,z]\mapsto[x,y]$ 给出的复\trManifold{}之间的\trMap{} $\mathbb{CP}^2\to\mathbb{CP}^1$，因为当 $x=y=0$ 时该\trMap{}没有定义。这对应于如下事实：由 $s\mapsto x$ 与 $t\mapsto y$ 给出的\trGradedRing{}的\trMap{} $\mathbb{C}[s,t]\to\mathbb{C}[x,y,z]$ 不满足假设 \eqref{eq:vphi-s-plus-empty}。

\textbf{7.4.B.（练习）}\quad 证明：若 $\phi\colon S_\bullet\to R_\bullet$ 满足 $\sqrt{(\phi(S_+))}=R_+$，则假设 \eqref{eq:vphi-s-plus-empty} 成立。（提示：见\cref{ex:empty-vanishing-equivalents}。）这个更为几何化的假设的代数表述，有时验证起来更为容易。

\textbf{7.4.C.（不重要的练习）}\quad 这道练习表明：不同的\trGradedRing{}的\trMap{}可以给出同一个\trScheme{}的\trMap。设 $R_\bullet=k[x,y,z]/(xz,yz,z^2)$、$S_\bullet=k[a,b,c]/(ac,bc,c^2)$，其中每个变量的次数都是 $1$。证明：$\Proj R_\bullet\cong\Proj S_\bullet\cong\mathbb{P}^1_k$。证明：由 $(a,b,c)\mapsto(x,y,z)$ 与 $(a,b,c)\mapsto(x,y,0)$ 给出的那两个\trMap{} $S_\bullet\to R_\bullet$，给出同一个（同）构 $\Proj R_\bullet\to\Proj S_\bullet$。（真正的缘由是：所有这些构造对「在有穷多个次数上发生的事」都不敏感。这一点日后将从多个途径说得精确，最近的便是练习~7.4.F。）

\textbf{7.4.3.}\ \emph{不重要的注记。}\quad \mbox{练习~15.2.H} 表明：即便在相当合理的情形下，也并非\trScheme{}的每个\trMorphism{} $\Proj R_\bullet\to\Proj S_\bullet$ 都来自\trGradedRing{}的\trMap{} $S_\bullet\to R_\bullet$。

\textbf{7.4.4.\ \trVeroneseSubring{}。}\quad 这里有一个有用的构造。设 $S_\bullet$ 是一个\trFinitelyGenerated\trGradedRing。把 $S_\bullet$ 的第 $n$ 个\textbf{\trVeroneseSubring{}}（$n$th Veronese subring）定义为 $S_{n\bullet}=\oplus_{j=0}^{\infty}S_{nj}$。（「旧次数」$n$ 即「新次数」$1$。）这一构造的几何解读，就是重要的\emph{\trVeroneseEmbedding{}}（Veronese embedding），将于 \S9.3.6 中讨论。

\textbf{7.4.D.（练习）}\quad 证明：\trGradedRing{}的\trMap{} $S_{n\bullet}\hookrightarrow S_\bullet$ 诱导一个\emph{\trIsomorphism} $\Proj S_\bullet\mathrel{\xrightarrow{\ \sim\ }}\Proj S_{n\bullet}$。（提示：若 $f\in S_+$ 是次数被 $n$ 整除的\trHomogeneous{}元素，把 $\Proj S_\bullet$ 上的 $D(f)$ 与 $\Proj S_{n\bullet}$ 上的 $D(f)$ 等同起来。为什么这样的\trDistinguishedOpen{}覆盖 $\Proj S_\bullet$？）

\textbf{7.4.E.（练习）}\quad 若 $S_\bullet$ 在次数 $1$ 处生成，证明 $S_{n\bullet}$ 也在次数 $1$ 处生成。（你或许愿意先考虑\trPolynomialRing{}的情形。）

\textbf{7.4.F.（练习）}\quad 证明：若 $R_\bullet$ 与 $S_\bullet$ 是同一个\trFinitelyGenerated\trGradedRing，仅在有穷多个非零次数上彼此不同（请把这话说精确！），则 $\Proj R_\bullet\cong\Proj S_\bullet$。

\textbf{7.4.G.（练习）}\quad 设 $S_\bullet$ 在 $S_0$ 上由 $f_1,\dots,f_n$ 生成。找出一个 $d$，使得 $S_{d\bullet}$ 在「新」次数 $1$（即「旧」次数 $d$）处是\trFinitelyGenerated{}的。（这出人意料地棘手，所以给一条提示。设有\trGenerator{} $x_1,\dots,x_n$，次数依次为 $d_1,\dots,d_n$。证明：任何次数至少为 $nd_1\cdots d_n$ 的\trMonomial{} $x_1^{\alpha_1}\cdots x_n^{\alpha_n}$，必有某个 $i$ 使得 $\alpha_i\ge(\prod_j d_j)/d_i$。再证明：第 $nd_1\cdots d_n$ 个\trVeroneseSubring{}由「新」次数 $1$ 处的元素生成。）

练习~7.4.G 与练习~7.4.D 结合起来表明：假定\trFinitelyGenerated\trGradedRing{}都在次数 $1$ 处生成，几乎无害——因为在\trRegrading{}之后（更精确地说：只保留次数为 $d$ 之倍数的那些项，再把次数除以 $d$），情形确实如此。这一点很便利，因为它意味着：借练习~7.4.D 之助，我们可以假定任何\trFinitelyGenerated\trGradedRing{}都在次数 $1$ 处生成。\mbox{练习~9.3.F} 日后将作为一个推论蕴含：我们可以把每个 $\Proj$ 嵌入某个\trProjectiveSpace{}之中。

\textbf{7.4.H.（较不重要的练习）}\quad 设 $S_\bullet$ 是一个\trFinitelyGenerated\trGradedRing。证明：$S_{n\bullet}$ 是一个\trFinitelyGenerated\trGradedRing。（一个可能的途径：用上一道练习，或某个类似的办法，证明存在某个 $N$ 使得 $S_{nN\bullet}$ 在次数 $1$ 处生成，从而\trGradedRing{} $S_{nN\bullet}$ 是\trFinitelyGenerated{}的。然后证明：对每个 $0<j<N$，$S_{nN\bullet+nj}$ 是 $S_{nN\bullet}$ 上的一个\trFinitelyGenerated{}的\trModule。）


\section{既约概形上的\trRationalMap{}}

不严格地说，「\trRationalMap{}」就是「几乎处处有定义的\trMorphism」，正如\trRationalFunction{}%
（定义~6.6.36、练习~5.2.I）乃是「几乎处处有定义的函数」的一个名号。我们日后将会看到：
在良好的情形下，正如有理函数那样，\trRationalMap{}在其有定义之处取值是唯一的（「既约到分离」%
定理~11.3.2，Reduced-to-Separated Theorem），并且拥有一个极大的「\trDomainOfDefinition{}」（\S11.3.3）。仅就本节而言，%
\emph{我们假定 $X$ 是\trReduced{}的}，尽管如此，我们仍会尽量小心，在所有定理中把这些假设%
明说出来。一个关键的例子将是\trIrreducible\trVariety（\S7.5.7），而\trRationalMap{}的语言%
也正是用在这种场合最为频繁。

\textbf{7.5.1.}\ \emph{\trhDefinition{}。}\quad 从一个\trReduced{}\trScheme{} $X$ 到一个%
\trScheme{} $Y$ 的\textbf{\trRationalMap{}}（rational map）$\pi$，记作 $\pi\colon X\dashrightarrow Y$，
是指如下数据：从一个\trDense\trOpenSubset{} $U\subset X$ 出发的一个\trMorphism{}%
$\alpha\colon U\to Y$，并配有如下的\trEquivalenceRelation：%
$(\alpha\colon U\to Y)\sim(\beta\colon V\to Y)$，若存在一个\trDense\trOpenSet{}%
$Z\subset U\cap V$ 使得 $\alpha|_Z=\beta|_Z$。（在 \S11.3.3 中，我们将把它改进为：在良好的%
情形下——即 $Y$ \trSeparated{}时——只须 $\alpha|_{U\cap V}=\beta|_{U\cap V}$。）人们常把%
「映射」一词当作「态射」的同义词使用，这相当合理；只不过\trRationalMap{}未必是一个映射。因此为了%
避免混淆，当你的意思是「\trRationalMap{}」时，切不可只说「映射」。

我们也将讨论\trScheme{} $S$ 上诸 $S$-\trScheme{}的\trRationalMap{}。定义照旧，只是现在 $X$ 与 $Y$ %
都是 $S$-\trScheme，而 $\alpha\colon U\to Y$ 是 $S$-\trScheme{}之间的一个\trMorphism。

\textbf{7.5.2.}\ \hard\ \emph{更一般的\trRationalMap{}。}\quad 正如有理函数那样，定义~7.5.1 %
也可以推广到 $X$ 不\trReduced{}的情形，做法要么照搬（沿用同一个名字「\trRationalMap{}」），要么%
采用一个在非\trReduced\trLocus{}上施加某种控制的版本（\textbf{拟态射}（pseudo-morphisms），
[Stacks, tag 01RX]）。我们将在 \S11.3 中看到，从\trReduced{}\trScheme{}到\trSeparated{}%
\trScheme{}的\trRationalMap{}表现得格外好，这正是人们通常在这一语境下考虑它们的原因。拟态射这一%
定义的用意，便是把这些结果推广到 $X$ 非\trReduced{}的情形。我们不会用到拟态射这个观念。

\textbf{7.5.3.}\quad \trRationalMap{}的一个显然例子是\trMorphism{}本身。另一个重要的例子是投影%
$\mathbb{P}^n_A\dashrightarrow\mathbb{P}^{n-1}_A$，它由 $[x_0,\dots,x_n]\mapsto[x_0,\dots,x_{n-1}]$ 给出。%
（确切地说，按定义~7.5.1 的意义，这凭什么是一个\trRationalMap{}？）

\textbf{7.5.A.（练习）}\quad 设 $X$ 是一个\trReduced{}\trScheme。请描述 $X\dashrightarrow%
\mathbb{A}^1$ 型的\trRationalMap{}应当如何与 $X$ 上的\trRationalFunction{}相认同。

\trRationalMap{} $\pi\colon X\dashrightarrow Y$ 称为\textbf{\trDominant{}}（dominant）的（在某些文献里%
亦作 dominating），若对其某一个代表 $U\to Y$——从而也就是对其每一个代表——其像在 $Y$ 中%
\trDense。一个\trMorphism{}称为\textbf{\trDominant{}态射}（dominant morphism，或 dominating %
morphism），若它作为\trRationalMap{}是\trDominant{}的。

\textbf{7.5.B.（练习）}\quad 证明：\trIntegral{}\trScheme{}之间的\trRationalMap{}%
$\pi\colon X\dashrightarrow Y$ 是\trDominant{}的，当且仅当 $\pi$ 把 $X$ 的\trGenericPoint{}送到%
$Y$ 的\trGenericPoint。

\textbf{7.5.4.}\ \emph{\trRationalMap{}的复合。}\quad 稍加思索你便会相信：可以把一个从\trReduced{}%
\trScheme{} $X$ 到\trReduced{}\trScheme{} $Y$ 的\trDominant{}映射 $\pi\colon X\dashrightarrow Y$，%
与一个\trRationalMap{} $\rho\colon Y\dashrightarrow Z$ 复合起来（并且复合是\trWellDefined{}的）。%
此外，若 $\rho$ 是\trDominant{}的，则复合 $\rho\circ\pi$ 也是\trDominant{}的。\trIntegral{}\trScheme{}连同%
其间的\trDominant{}\trRationalMap{}构成一个在几何上饶有趣味的\trCategory。

\textbf{7.5.C.（容易的练习）}\quad 证明：\trIntegral{}\trScheme{}之间的\trDominant{}\trRationalMap{}按%
相反的方向给出\trFunctionField{}之间的\trMorphism。

在「适当经典的情形」（\trIntegral\trFiniteType{}的 $k$-\trScheme——特别地，还有将在%
\S11.2.9 中定义的\trIrreducible\trVariety）中，这件事是可逆的：\trDominant{}\trRationalMap{}对应于%
反方向的函数域的包含。我们将在下文 \S7.5.7 中把这一点说精确。但是，「函数域之间的%
\trMorphism{}总能给出\trDominant{}\trRationalMap{}」并不成立，甚至未必能给出\trRationalMap{}。例如，$\Spec k[x]$ %
与 $\Spec k(x)$ 具有同一个\trFunctionField{} $k(x)$，却不存在与之相应的 $k$-\trScheme{}%
的\trRationalMap{} $\Spec k[x]\dashrightarrow\Spec k(x)$。理由：这样的\trRationalMap{}应当对应于一个从%
$\Spec k[x]$ 的某个\trOpenSubset{} $U$（比如说 $\Spec k[x,1/f(x)]$）到 $\Spec k(x)$ %
的\trMorphism。但对任何一个 $f(x)$ 而言，（把 $k$ 逐点恒同地送到 $k$、把 $x$ 送到 $x$ %
的）环映射 $k(x)\to k[x,1/f(x)]$ 都不存在。

（如果你还想要更多的证据，来说明这个用拓扑方式定义的\trDominant{}观念同时也是代数的，你可以%
证明：若 $\phi\colon A\to B$ 是一个\trRingMorphism，则相应的\trMorphism{}%
$\Spec B\to\Spec A$ 是\trDominant{}的，当且仅当 $\phi$ 的\trKernel{}含于 $A$ 的\trNilradical{}%
之中。）

\textbf{7.5.5.}\ \emph{\trhDefinition{}。}\quad \trIntegral{}\trScheme{}之间的\trRationalMap{}%
$\pi\colon X\dashrightarrow Y$ 称为\textbf{\trBirational{}}（birational）的，若它是\trDominant{}的，并且%
存在另一个同样是\trDominant{}的\trRationalMap{}（一个「有理逆」）$\psi$，使得 $\pi\circ\psi$（在同一个%
\trEquivalenceClass{}中）就是 $Y$ 上的恒同，且 $\psi\circ\pi$（在同一个\trEquivalenceClass{}%
中）就是 $X$ 上的恒同。这就是「\trIntegral{}\trScheme{}连同其间的\trDominant{}\trRationalMap{}」这一%
\trCategory{}中的\trIsomorphism{}观念。（在可微范畴里，这并不特别有趣。）一个\trMorphism{}%
若作为\trRationalMap{}是\trBirational{}的，便称为\textbf{\trBirational{}态射}（birational morphism）。（注意：%
\trBirational{}态射的「逆」可以只是一个\trRationalMap{}，而不是\trMorphism。）

若存在一个双\trRationalMap{} $X\dashrightarrow Y$，我们便说 $X$ 与 $Y$（相互）\trBirational{}。%
若 $X$ 与 $Y$ 都\trIrreducible，则双\trRationalMap{}诱导\trFunctionField{}之间的\trIsomorphism。%
「对\trIntegral\trFiniteType{}的 $k$-\trScheme，函数域之间的映射对应于反方向的\trRationalMap{}」%
这一事实（将于命题~7.5.8 中证明）表明：在\trIntegral\trFiniteType{}的 $k$-\trScheme{}%
之间，凡诱导\trFunctionField{}的\trIsomorphism{}的映射必是\trBirational{}的。一个%
\trIntegral\trFiniteType{}的 $k$-\trScheme{}若是\trBirational{}于某个 $\mathbb{A}^n_k$，便称它是%
\textbf{有理}（rational）的。

\textbf{7.5.6.\ \trhProposition。}\ ——\ \emph{设 $X$ 与 $Y$ 是既\trIrreducible{}又%
\trReduced{}的\trScheme（据练习~5.2.G，即\trIntegral{}的概形）。则 $X$ 与 $Y$ \trBirational{}，%
当且仅当存在 $X$ 的一个\trDense\trOpenSubscheme{} $U$ 与 $Y$ 的一个\trDense%
\trOpenSubscheme{} $V$，使得 $U\cong V$。}

（我们加上「整」的假设只是为了简单。配以足够的细心与修改过的定义，这些假设是可以去掉的。）

命题~7.5.6 告诉你该如何看待双\trRationalMap{}。正如\trRationalMap{}是「大半处有定义的函数」，两个\trBirational{}%
的\trReduced{}\trScheme{}便是「大半处同构」的。例如，一个\trReduced\trFiniteType{}的%
$k$-\trScheme（比如一个\trAffine{} $k$-\trVariety）是\trRational{}的，当且仅当它有一个%
\trIsomorphic{}于 $\mathbb{A}^n_k$ 的某个\trOpenSubscheme{}的\trDense\trOpenSubscheme。

\emph{证明。}\ 「当」方向是显然的，故我们证明「仅当」方向。我们基本上就是顺着鼻子走，%
把手头给的条件全部用上（见图~\ref{fig:birational-proof-outline}）。

\begin{figure}[htbp]
\centering
{\scriptsize
\begin{tikzcd}[column sep=1.35em, row sep=2.55em]
f & g\circ f & g\circ f=\id\ \text{作为：} & & & g & f\circ g & f\circ g=\id\ \text{作为：} \\
\text{\trRationalMap{}} & \text{\trRationalMap{}} & \text{\trRationalMap{}} &
X \ar[r,<->,dashed,"\sim"] & Y &
\text{\trRationalMap{}} & \text{\trRationalMap{}} & \text{\trRationalMap{}} \\
\text{态射} & \text{”} & \text{”} &
X_1 \ar[u,hook'] \ar[ur] & Y_1 \ar[u,hook'] \ar[ul] &
\text{态射} & \text{”} & \text{”} \\
\text{”} & \text{态射} & \text{”} &
X_2 \ar[u,hook'] \ar[ur] & & & & \\
\text{”} & \text{”} & \text{态射} &
X_3 \ar[u,hook'] & Y_2 \ar[uu,hook'] \ar[l] &
\text{”} & \text{态射} & \text{”} \\
\text{”} & \text{”} & \text{态射} &
X_4 \ar[u,hook'] \ar[r,<->,"\sim"] & Y_3 \ar[u,hook'] &
\text{”} & \text{”} & \text{态射}
\end{tikzcd}}
\caption{命题~7.5.6 的证明概观。（依原书手绘图重绘。）}
\label{fig:birational-proof-outline}
\end{figure}

设 $F\colon X\dashrightarrow Y$ 与 $G\colon Y\dashrightarrow X$ 是一对互逆的\trRationalMap{}。%
选取\trRationalMap{} $F$ 与 $G$ 各自的代表\trMorphism{} $f\colon X_1\to Y$（其中 $X_1\subset X$）%
与 $g\colon Y_1\to X$（其中 $Y_1\subset Y$）。令 $X_2=(f|_{X_1})^{-1}(Y_1)\subset X_1$，%
则 $g\circ f$ 是一个从 $X_2$ 到 $X$ 的\trMorphism，它作为\trRationalMap{}是恒同。于是存在一个%
\trDense\trOpenSubset{} $X_3\subset X_2$，使得\trMorphism{} $g\circ f\colon X_3\to X$ %
是到 $X_3$ 上的恒同\trMorphism（或者更确切地说，它是\trOpenEmbedding{}%
$X_3\hookrightarrow X$）。

类似地，令 $Y_2=(g|_{Y_1})^{-1}(X_3)\subset Y_1$，于是 $f\circ g$ 是一个从 $Y_2$ 到 $Y$ %
的\trMorphism，它作为\trRationalMap{}是恒同。再令 $Y_3\subset Y_2$ 为一个\trDense\trOpenSubset，%
使得 $f\circ g\colon Y_3\to Y$ 是包含映射（即「恒同」）。

最后，令 $X_4=(f|_{X_3})^{-1}(Y_3)\subset X_3$，则 $(g\circ f)|_{X_4}$（经由 $Y_3$）%
是 $X_4$ 上的恒同\trMorphism，而 $(f\circ g)|_{Y_3}$（经由 $X_4$）是 $Y_3$ 上的恒同%
\trMorphism，于是我们便找到了这个开集之间的\trIsomorphism，它正是我们这个双\trRationalMap{}的%
一个「代表」。\hfill$\square$

\textbf{7.5.7.\ \trIrreducible\trVariety{}的\trRationalMap{}。}

\textbf{7.5.8.\ \trhProposition。}\ ——\ \emph{设 $X$ 是一个\trIntegral{}的 $k$-\trScheme，%
$Y$ 是一个\trIntegral\trFiniteType{}的 $k$-\trScheme，并给定一个保持 $k$ 的%
\trFunctionField{}扩张 $\phi^{\sharp}\colon K(Y)\hookrightarrow K(X)$。则存在 $k$-\trScheme{}%
的一个\trDominant{}\trRationalMap{} $\phi\colon X\dashrightarrow Y$ 诱导 $\phi^{\sharp}$。}

\emph{证明。}\ 用一个\trOpenSubset{}替换 $Y$ 之后，不妨设 $Y$ 是\trAffine{}的，比如说%
$\Spec B$，其中 $B$ 由有限多个元素 $y_1,\dots,y_n$ 在 $k$ 上生成。由于我们只须在 $X$ %
的一个\trOpenSubset{}上定义 $\phi$，同样不妨设 $X=\Spec A$ 是\trAffine{}的。于是%
$\phi^{\sharp}$ 给出一个包含 $\phi^{\sharp}\colon B\hookrightarrow K(A)$。把 $\phi^{\sharp}(y_i)$ %
写成 $f_i/g_i$（其中 $f_i,g_i\in A$），并令 $g:=\prod_i g_i$。于是 $\phi^{\sharp}$ 进而%
诱导一个包含 $B\hookrightarrow A_g$。从而 $\phi\colon\Spec A_g\to\Spec B$ 诱导%
$\phi^{\sharp}$。\trMorphism{} $\phi$ 是\trDominant{}的，因为零\trIdeal{}在包含 $B\hookrightarrow A_g$ %
之下的逆像仍是零\trIdeal，于是 $\phi$ 把 $X$ 的\trGenericPoint{}送到 $Y$ 的%
\trGenericPoint。\hfill$\square$

\textbf{7.5.D.（练习）}\quad 设 $K$ 是 $k$ 的一个\trFinitelyGenerated\trFieldExtension。%
（回想：$k$ 上的域扩张 $K$ 称为\textbf{有限生成}的，若 $K$ 中存在一个有限的「生成集」%
$\{x_1,\dots,x_n\}$，使得 $K$ 的每个元素都能写成系数属于 $k$ 的、以 $x_1,\dots,x_n$ 为%
变元的\trRationalFunction。）证明：存在以 $K$ 为\trFunctionField{}的\trIrreducible%
\trAffine{} $k$-\trVariety。（提示：考虑由 $t_i\mapsto x_i$ 给出的映射%
$k[t_1,\dots,t_n]\to K$，证明其\trKernel{}是一个\trPrimeIdeal{} $\mathfrak{p}$，且%
$k[t_1,\dots,t_n]/\mathfrak{p}$ 的\trFractionField{}为 $K$。几何解读：考虑由环映射%
$t_i\mapsto x_i$ 给出的映射 $\Spec K\to\Spec k[t_1,\dots,t_n]$，并取这个单点像的%
\trClosure。）

\textbf{7.5.E.（练习）}\quad 请描述下列诸范畴之间的\trEquivalenceOfCategories。

(a)~以「\trIntegral\trAffine{} $k$-\trVariety」为\trObject、「在 $k$ 上定义的\trDominant{}有理%
映射」为\trMorphism{}的范畴；以及

(b)~以「$k$ 的\trFinitelyGenerated\trFieldExtension」为\trObject、「在 $k$ 上延拓恒同%
的包含」为\trMorphism{}的（「箭头掉头」的）\trOppositeCategory。

（待我们在 \S11.2.9 中定义了一般的\trVariety{}之后，还可以添上：(c)~以「\trIntegral{} %
$k$-\trVariety」为\trObject、「在 $k$ 上定义的\trDominant{}\trRationalMap{}」为\trMorphism{}的范畴。）

特别地，一个\trIntegral\trAffine{} $k$-\trVariety{} $X$ 是\trRational{}的，若其\trFunctionField{}%
$K(X)$ 是 $k$ 的一个\textbf{纯超越扩张}（purely transcendental extension），也就是说%
$K(X)\cong k(x_1,\dots,x_n)$（对某个 $n$）。（这还需要说得更精确些：由 $X\to\Spec k$ %
诱导的映射 $k\leftarrow K(X)$，应当在这个\trIsomorphism{}之下与「显然的」映射%
$k\leftarrow k(x_1,\dots,x_n)$ 相合。）

\textbf{7.5.9.\ \trRationalMap{}的更多例子。}

这些例子里反复出现的一个主题是：到\trProjectiveScheme{}的\trRationalMap{}，其\trDomainOfDefinition{}%
可以越过余维数为 $1$ 的\trRegular{}点向上延拓。等我们讨论\trCurve{}时，将在「曲线到射影%
扩张定理」（Curve-to-Projective Extension Theorem）~15.3.1 中把这一点说精确。

\begin{figure}[htbp]
\centering
\begin{tikzpicture}[scale=1.15]
  \draw[thick] (0,0) circle (1);
  \draw[->] (-1.7,0) -- (1.85,0) node[right] {$x$};
  \draw[->] (0,-1.7) -- (0,1.85) node[right] {$y$};
  \fill (1,0) circle (1.7pt) node[above right] {$p$};
  \fill (-0.45,0.89) circle (1.7pt) node[above left] {$q$};
  \draw[thick] (-1.04,1.25) -- (1.65,-0.4);
  \node at (-1.2,-0.8) {$C$};
  \node at (1.5,-0.68) {斜率 $m$};
\end{tikzpicture}
\caption{用几何寻找本原勾股三元组。（依原书手绘图重绘。）}
\label{fig:pythagorean-stereographic}
\end{figure}

第一个例子是\textbf{勾股三元组}（Pythagorean triples）的经典公式，以及它经由「球极投影」%
（stereographic projection）的推导。设你正在寻找由 $x^2+y^2=1$ 给出的圆 $C$ %
（图~\ref{fig:pythagorean-stereographic}）上的\trRationalPoint。一个\trRationalPoint{}%
是 $p=(1,0)$。若 $q$ 是另一个有理点，则 $pq$ 是一条斜率为有理数（非无穷）的直线。这给出%
一个从\trConic{} $C$（现在解读为 $\Spec\mathbb{Q}[x,y]/(x^2+y^2-1)$）到 $\mathbb{A}^1_{\mathbb{Q}}$ %
的\trRationalMap{}，由 $(x,y)\mapsto y/(x-1)$ 给出。（刚才悄悄发生了一件微妙的事情：我们本来%
谈的是圆上的 $\mathbb{Q}$-点，结果却得到了概形之间的一个\trRationalMap{}。）反过来，给定过 $p$、%
斜率为 $m$（$m$ 为有理数）的一条直线，我们可以通过解方程组 $y=m(x-1)$ 与 $x^2+y^2=1$ %
把 $q$ 找回来。把第一个方程代入第二个，便得到一个关于 $x$ 的二次方程。我们知道 $x=1$ %
必是一个解（因为该直线与圆相交于 $(x,y)=(1,0)$），所以我们指望能把这个因子提出来，%
找出另一个因子。这确实行得通：
\[
\begin{gathered}
x^2+\bigl(m(x-1)\bigr)^2=1\\
\implies\quad (m^2+1)x^2+(-2m^2)x+(m^2-1)=0\\
\implies\quad (x-1)\bigl((m^2+1)x-(m^2-1)\bigr)=0.
\end{gathered}
\]
另一个解是 $x=(m^2-1)/(m^2+1)$，由此得 $y=-2m/(m^2+1)$。于是我们得到\trConic{} $C$ %
与以 $m$ 为坐标的 $\mathbb{A}^1$ 之间的一个双\trRationalMap{}，由 $f\colon(x,y)\mapsto y/(x-1)$ %
给出（它在 $x\neq 1$ 处有定义），其逆\trRationalMap{}由 $m\mapsto\bigl((m^2-1)/(m^2+1),\,-2m/(m^2+1)\bigr)$ %
给出（它在 $m^2+1=0$ 之外有定义）。

我们可以经由「包含」$\mathbb{A}^1_{\mathbb{Q}}\to\mathbb{P}^1_{\mathbb{Q}}$（日后我们将%
称之为一个\trOpenEmbedding）把它扩充为一个\trRationalMap{} $C\dashrightarrow\mathbb{P}^1_{\mathbb{Q}}$。%
此时 $f$ 由 $(x,y)\mapsto[y,\,x-1]$ 给出。于是我们有了一个有趣的问题：$f$ 的%
\trDomainOfDefinition{}是什么？看上去它处处有定义，唯独在 $y=x-1=0$ 之处——也就是除 $p$ %
之外——没有定义。但事实上它可以越过 $p$ 延拓！注意：$(x,y)\mapsto[x+1,\,-y]$（其中%
$(x,y)\neq(-1,0)$）在二者公共的\trDomainOfDefinition{}上与 $f$ 相合，因为%
$[x+1,\,-y]=[y,\,x-1]$。于是这个\trRationalMap{}能延拓得比我们起初所想的更远。这将是曲线到射影%
扩张定理~15.3.1 的一个特例。

\textbf{7.5.F.（练习）}\quad 利用上述办法，找出一个给出全体勾股三元组的「公式」。

\textbf{7.5.G.（练习）}\quad 证明：对任何\trCharacteristic{}不为 $2$ 的\trField{} $k$，%
$\mathbb{P}^2_k$ 中的\trConic{} $x^2+y^2=z^2$ 都\trIsomorphic{}于 $\mathbb{P}^1_k$。%
（如有需要，「$x^2+y^2=z^2$ 在 $\mathbb{P}^2_k$ 中」的定义可参\cref{ex:vanishing-scheme-definition}。）

\textbf{7.5.10.}\quad 事实上，$\mathbb{P}^2_k$ 中任何带有一个秩为 $3$ 的 $k$-\trValuedPoint{}%
（即\trResidueField{}为 $k$ 的一个点）的\trConic{}都\trIsomorphic{}于 $\mathbb{P}^1_k$ %
（要在基变换到 $\overline{k}$ 之后「秩」才有意义；参练习~5.4.J）。（带有 $k$-\trValuedPoint{}%
这一假设当然是必要的：$k=\mathbb{R}$ 上的 $x^2+y^2+z^2=0$ 便是一个不\trIsomorphic{}于%
$\mathbb{P}^1_k$ 的\trConic。）

\textbf{7.5.H.（练习）}\quad 通过找一个到 $\mathbb{A}^1_{\mathbb{Q}}$ 的双\trRationalMap{}——%
照搬处理\trConic{}时行之有效的那一套——求出 $y^2=x^3+x^2$ 的全部有理解。提示：你该从%
哪个点出发作投影？（在练习~19.10.I 中我们将会看到：这些点基本上构成一个\trGroup，而且%
这是一条退化的\trEllipticCurve。）

你将得到一个到 $\mathbb{P}^1_{\mathbb{Q}}$ 的\trRationalMap{}，它在「结点」（node）$x=y=0$ 之上%
没有定义，而且不能越过这个余维数 $1$ 的集合延拓。这恰好展示了我们日后的结果——讲述如何%
把\trRationalMap{}越过余维数 $1$ 的集合延拓到射影空间的曲线到射影扩张定理~15.3.1——的局限：%
余维数 $1$ 的集合必须是\trRegular{}的。

\textbf{7.5.I.（练习）}\quad 用类似的想法，找出从\trQuadricSurface{}%
$Q=\{x^2+y^2=w^2+z^2\}\subset\mathbb{P}^3_{\mathbb{Q}}$ 到 $\mathbb{P}^2_{\mathbb{Q}}$ %
的一个双\trRationalMap{}。再用它求出 $Q$ 上全部的\trRationalPoint。（这展示了解\textbf{丢番图%
方程}（Diophantine equations）的一种好办法。你将会找到 $Q$ 的一个\trIsomorphic{}于%
$\mathbb{P}^2$ 的某个\trDense\trOpenSubset{}的\trDense\trOpenSubset，在那里你可以%
轻轻松松求出全部\trRationalPoint。$Q$ 还会有一个\trClosedSubset，使得\trRationalMap{}在其上%
没有定义、或者不是\trIsomorphism，但你可以就事论事（ad hoc）地处置这个子集。）

\textbf{7.5.J.（练习：\trCremona{}（Cremona）变换——一个有用的经典构造）}\quad 考虑由%
$[x,y,z]\mapsto[1/x,\,1/y,\,1/z]$ 给出的\trRationalMap{} $\mathbb{P}^2_k\dashrightarrow\mathbb{P}^2_k$。%
它的\trDomainOfDefinition{}是什么？（它比 $xyz\neq 0$ 的\trLocus{}要大！）你将会观察到：%
你可以把它越过「余维数 $1$ 的集合」延拓（姑且不管我们迄今还不知道余维数是什么意思）。%
这又一次预示了曲线到射影扩张定理~15.3.1。

\textbf{7.5.11.}\ \hard\ \textbf{非\trRational{}的复\trCurve{}（有趣但非必要）。}\quad 我们现在来%
描述两个不存在从 $\mathbb{P}^1_{\mathbb{C}}$ 出发的非常值\trRationalMap{}的\trCurve{} $C$ 的例子。%
（诚然，我们迄今还不知道「\trCurve」是什么意思，但无妨。）两个证明用的都是\trFermat{}（Fermat）%
的\textbf{无穷递降}（infinite descent）法。这些结果可以（日后你将能借助定理~16.3.3 来%
核验）解读为这样一个事实：这些曲线没有「非平凡」的 $\mathbb{C}(t)$-\trValuedPoint——%
我们的意思是，任何 $\mathbb{C}(t)$-\trValuedPoint{}都暗中是一个 $\mathbb{C}$-\trValuedPoint。%
（你或许会注意到：若把同样的例子中的 $\mathbb{C}(t)$ 换成 $\mathbb{Q}$（并且 $C$ 是%
$\mathbb{Q}$ 而非 $\mathbb{C}$ 上的曲线），你就得到了数论与几何中的两个基本问题。%
练习~7.5.M 的类比是\trEllipticCurve{}上\trRationalPoint{}的问题，而你也许会意识到，%
练习~7.5.K 的类比甚至还要更出名。此外，练习~7.5.M(a) 在算术上的类比是「四平方数定理」%
（Four Squares Theorem）——不存在四个构成等差数列的整数平方——由\trFermat{}首次陈述。）

这些例子将让你得以一窥：数域上的事实为何、又如何常常有曲线的\trFunctionField{}上的%
事实与之平行。这种平行性是本学科中一个反复出现的深刻主题。

\textbf{7.5.K.（练习）}\quad 设 $n>2$，证明：从 $\mathbb{P}^1_{\mathbb{C}}$ 到%
$\mathbb{P}^2_{\mathbb{C}}$ 中的「\trFermat{}曲线」$x^n+y^n=z^n$ 不存在\trDominant{}\trRationalMap{}。提示：%
把它化归为证明：$f(t)^n+g(t)^n=h(t)^n$ 不存在「非常值」的解 $\bigl(f(t),g(t),h(t)\bigr)$，%
这里 $f(t)$、$g(t)$、$h(t)$ 是 $t$ 的\trRationalFunction（也就是 $\mathbb{C}(t)$ 的%
元素）。通过清去分母，再把它化归为证明：当 $f(t)$、$g(t)$、$h(t)$ 为互素多项式时，%
无非常值的解。为此，假设有解存在，并取其中正次数最低者加以考虑。然后利用 $\mathbb{C}[t]$ %
是\trUFD{}这一事实，以及等式 $h(t)^n-g(t)^n=\prod_{i=1}^{n}\bigl(h(t)-\zeta^i g(t)\bigr)$ %
（其中 $\zeta$ 是本原 $n$ 次单位根）。论证每个 $h(t)-\zeta^i g(t)$ 都是一个 $n$ 次幂。%
然后用
\[
\bigl(h(t)-g(t)\bigr)+\alpha\bigl(h(t)-\zeta g(t)\bigr)=\beta\bigl(h(t)-\zeta^2 g(t)\bigr)
\]
（其中 $\alpha$、$\beta$ 适当选取）来得到一个次数更低的解。（当 $n=2$ 时，这个论证%
为何失效？）

\textbf{7.5.L.（练习）}\quad 给出两条\trSmooth{}复\trCurve{} $X$ 与 $Y$，使得 $X$ 的任何%
非空\trOpenSubset{}都不\trIsomorphic{}于 $Y$ 的任何非空\trOpenSubset。（请试着不要为%
我们尚未定义「\trSmoothness」这件事而烦恼。）提示：练习~7.5.K。

\textbf{7.5.M.（练习）}\quad 设 $a$、$b$、$c$ 是互不相同的复数。请按下列步骤证明：若%
$x(t)$ 与 $y(t)$ 是 $t$ 的两个\trRationalFunction，使得
\begin{equation*}\tag{7.5.11.1}\label{eq:y-squared-cubic-product}
y(t)^2=\bigl(x(t)-a\bigr)\bigl(x(t)-b\bigr)\bigl(x(t)-c\bigr),
\end{equation*}
则 $x(t)$ 与 $y(t)$ 都是常数（即 $x(t),y(t)\in\mathbb{C}$）。（这里的 $\mathbb{C}$ 可以%
换成任何\trCharacteristic{}不为 $2$ 的\trField{} $K$；若 $K$ 不\trAlgClosed，则还需%
稍加小心。）

(a)~设 $P,Q\in\mathbb{C}[t]$ 是互素的多项式，使得它们有四个线性组合是完全平方，且这些%
完全平方中的任何两个都互非常数倍。证明：$P$ 与 $Q$ 都是常数（即 $P,Q\in\mathbb{C}$）。%
提示：通过给 $P$、$Q$ 改名，证明你可以假定这些完全平方恰是 $P$、$Q$、$P-Q$、$P-\lambda Q$ %
（对某个 $\lambda\in\mathbb{C}$）。定义 $u$ 与 $v$ 分别为 $P$ 与 $Q$ 的平方根。证明：%
$u-v$、$u+v$、$u-\sqrt{\lambda}\,v$、$u+\sqrt{\lambda}\,v$ 都是完全平方，且 $u$ 与 $v$ %
互素。若 $P$ 与 $Q$ 不都是常数，注意 $0<\max(\deg u,\deg v)<\max(\deg P,\deg Q)$。%
从一开始就假定 $P$ 与 $Q$ 被选作反例时已使 $\max(\deg P,\deg Q)$ 达到最小，由此导出%
矛盾。（题外话：三个互不相同的线性组合完全可能是完全平方。这类例子本质上对应于%
$\mathbb{C}(t)$ 中的本原勾股三元组——你能看出是怎么回事吗？）

(b)~设 $(x,y)=(p/q,\,r/s)$ 是 \eqref{eq:y-squared-cubic-product} 的一个解，其中%
$p,q,r,s\in\mathbb{C}[t]$，且 $p/q$ 与 $r/s$ 已写成最简形式。清去分母，证明%
$r^2q^3=s^2(p-aq)(p-bq)(p-cq)$。证明 $s^2\mid q^3$ 且 $q^3\mid s^2$，从而对某个%
$\delta\in\mathbb{C}$ 有 $s^2=\delta q^3$。由 $r^2=\delta(p-aq)(p-bq)(p-cq)$，证明%
$(p-aq)$、$(p-bq)$、$(p-cq)$ 都是完全平方。再证明 $q$ 也是完全平方，然后援引 (a)。

练习~21.4.H 给出了处理练习~7.5.K 与练习~7.5.M 的一种好得多的几何途径。


\section{\hard\ 可表函子与群概形}

\textbf{7.6.1.\ 到 $\mathbb{A}^1$ 的映射对应于函数。}\quad 若 $X$ 是一个\trScheme，那么到 $\mathbb{A}^1$ 的映射与\trStructureSheaf{}的\trGlobalSections{}之间存在一个双射：由必作练习~7.3.G，映射 $\pi\colon X\to\mathbb{A}^1_{\mathbb{Z}}$ 对应于\trRingMorphism{} $\pi^{\sharp}\colon\mathbb{Z}[t]\to\Gamma(X,\OO_X)$，而 $\pi^{\sharp}(t)$ 就是 $X$ 上的一个函数；这件事是可逆的。

在一种不严格的意味上，这个映射是非常自然的：你甚至可以在脑海里把这个到 $\mathbb{A}^1$ 的映射想象成\emph{由}那个函数所\emph{给出}。（作个类比：\trManifold{}上的函数就是一个到 $\mathbb{R}$ 的映射。）但它还有一种更精确的意味下的自然：这个双射关于 $X$ 是\trFunctoriality{}的。我们将细细地把这个例子琢磨个够，并且将会看到：它把我们引向两个重要而不简单的观念——\trRepresentableFunctor{}与\trGroupScheme。

\textbf{7.6.A.（容易的练习）}\quad 设 $X$ 是一个 $\mathbb{C}$-\trScheme。请验证：$\mathbb{C}$-\trScheme{}的\trCategory{}中到 $\mathbb{A}^1_{\mathbb{C}}$ 的映射，与 $X$ 上的函数之间，存在一个自然的双射。（这里的\trBaseRing{} $\mathbb{C}$ 可以换成任何\trRing{} $A$。）

这一解读还可以推广到\trRationalMap，如下。

\textbf{7.6.B.（无关紧要的练习）}\quad 请把整\trScheme{}上的\trRationalFunction（\mbox{练习~5.2.I}）解读为到 $\mathbb{A}^1_{\mathbb{Z}}$ 的\trRationalMap。

（如果你愿意，还可以把你的论证推广到局部\trNoetherian{}\trScheme{}上的\trRationalMap；见定义~6.6.36。）

\textbf{7.6.2.\ \trRepresentableFunctor。}\quad 我们把 \S7.6.1 的那个双射换一种说法重述如下。我们有从 $\mathit{Sch}$ 到 $\catSets$ 的两个不同的\trContravariantFunctor：到 $\mathbb{A}^1$ 的映射（即 $\mathrm{H}\colon X\mapsto\Mor(X,\mathbb{A}^1_{\mathbb{Z}})$），以及 $X$ 上的函数（$\mathrm{F}\colon X\mapsto\Gamma(X,\OO_X)$）。这个双射的「自然性」——也就是它关于 $X$ 的\trFunctoriality——恰好就是这样一个陈述：该双射给出诸函子间的一个\trNaturalIsomorphism（\S1.1.21）。也就是说，给定任何 $\pi\colon X\to X'$，图
\[
\begin{tikzcd}[row sep=3.2em]
\mathrm{H}(X') \ar[r, "\mathrm{H}(\pi)"] \ar[d] & \mathrm{H}(X) \ar[d] \\
\mathrm{F}(X') \ar[r, "\mathrm{F}(\pi)"'] & \mathrm{F}(X)
\end{tikzcd}
\]
（其中竖直的映射就是 \S7.6.1 中给出的那两个双射）交换。

更一般地，若 $Y$ 是\trCategory{} $\scr{C}$ 的一个元素（我们在意的特例是 $\scr{C}=\mathit{Sch}$），回想由 $\mathrm{h}_Y(X)=\Mor(X,Y)$ 定义的\trContravariantFunctor{} $\mathrm{h}_Y\colon\scr{C}\to\catSets$（例~1.1.20）。我们说：一个从 $\scr{C}$ 到 $\catSets$ 的\trContravariantFunctor{}是\textbf{被 $Y$ 所\trRepresentation{}的}（represented by $Y$），若它与函子 $\mathrm{h}_Y$ \trNaturalIsomorphism；又说它是\textbf{\trRepresentable{}的}（representable），若它被\emph{某个} $Y$ 所\trRepresentation。

\S7.6.1 的那个双射，如今可以重述为：\emph{整体截面函子被 $\mathbb{A}^1$ 所\trRepresentation。}

\textbf{7.6.C.（重要而容易的练习：表示对象在相差唯一\trIsomorphism{}的意义下唯一）}\quad 证明：若一个\trContravariantFunctor{} $F$ 既被 $Y$ 所\trRepresentation{}又被 $Z$ 所\trRepresentation，那么便有一个唯一的\trIsomorphism{} $Y\mathrel{\xrightarrow{\ \sim\ }}Z$，它由函子的\trNaturalIsomorphism{} $\mathrm{h}_Y\mathrel{\xrightarrow{\ \sim\ }}\mathrm{h}_Z$ 所诱导。提示：这是 \S1.2 的泛性质论证的一个版本：我们再一次凭着映入某对象的那些映射认出该对象（在相差唯一\trIsomorphism{}的意义下）。这道练习本质上就是\cref{ex:general-yoneda-representables}(b)。

（在此框架下，这件事可以顺理成章地推广到\trYonedaLemma，即\cref{ex:yoneda-lemma-schemes}。你现在就去试上一试，亦无不可。）

其实你早已暗中见过这个观念：你可以把\trCategory{}中\trProduct{}与\trFiberProduct{}的存在性，解读为\trRepresentableFunctor{}的例子。（你或许愿意琢磨一下：$\mathrm{h}_{Y\times Z}\mathrel{\xrightarrow{\ \sim\ }}\mathrm{h}_Y\times\mathrm{h}_Z$ 的一个\trNaturalIsomorphism{}如何诱导出投影映射 $Y\times Z\to Y$ 与 $Y\times Z\to Z$。）

\textbf{7.6.D.（热身练习）}\quad 设 $F$ 是\trContravariantFunctor{} $\mathit{Sch}\to\catSets$，定义为：对一切概形 $X$，$F(X)=\{\text{Grothendieck}\}$。证明：$F$ 是\trRepresentable{}的。（它被谁所\trRepresentation？）\footnote{译注：笑点在于——「对一切 $X$ 恒取同一个单点集 $\{\text{Grothendieck}\}$」的函子，恰被 $\Spec\mathbb{Z}$ 表示：单点集的「万有元素」正是那位什么都研究的人。}

\textbf{7.6.E.（练习）}\quad 本练习中，$\mathbb{Z}$ 可以换成任何\trRing。
\begin{enumerate}[label=(\alph*)]
  \item \emph{（仿射 $n$-空间\trRepresentation{}「$n$ 个函数」的函子。）}证明：从（$\mathbb{Z}$-）概形到 $\catSets$ 的\trContravariantFunctor
  \[
  X\longmapsto\{(f_1,\dots,f_n)\ :\ f_i\in\Gamma(X,\OO_X)\}
  \]
  被 $\mathbb{A}^n_{\mathbb{Z}}$ 所\trRepresentation。证明：$\mathbb{A}^1_{\mathbb{Z}}\times_{\mathbb{Z}}\mathbb{A}^1_{\mathbb{Z}}\cong\mathbb{A}^2_{\mathbb{Z}}$；也就是说，$\mathbb{A}^2$ 满足 $\mathbb{A}^1\times\mathbb{A}^1$ 的函子式刻画。（你无疑立刻就能顺手证明：$\prod\mathbb{A}^{m_i}_{\mathbb{Z}}\mathrel{\overset{\sim}{\longleftrightarrow}}\mathbb{A}^{\sum m_i}_{\mathbb{Z}}$。）
  \item \emph{（可逆函数的函子是可表的。）}证明：从（$\mathbb{Z}$-）概形到 $\catSets$、把 $X$ 送到 $X$ 上的\trInvertible{}函数的那个\trContravariantFunctor，被 $\Spec\mathbb{Z}[t,t^{-1}]$ 所\trRepresentation。
\end{enumerate}

\textbf{7.6.3.}\ \emph{\trhDefinition{}。}\quad \mbox{练习~7.6.E}(b) 中定义的那个概形，称为\textbf{\trMultiplicativeGroup{}}（multiplicative group）（或\textbf{\trMultiplicativeGroup{}概形}（multiplicative group scheme））$\mathbb{G}_m$。「域 $k$ 上的 $\mathbb{G}_m$」（「$k$ 上的\trMultiplicativeGroup{}」）指 $\Spec k[t,t^{-1}]$，并带着同样的群运算。更好的说法是：它在 $k$-\trScheme{}的\trCategory{}中\trRepresentation{}由可逆函数组成的群。我们同样可以在任意\trRing{}上、乃至任意\trScheme{}上定义 $\mathbb{G}_m$。

\textbf{7.6.F.（较不重要的练习）}\quad 固定一个\trRing{} $A$。考虑从\trLocallyRingedSpace{}的\trCategory{}到 $\catSets$ 的函子 $\mathrm{H}$，由 $\mathrm{H}(X)=\{A\to\Gamma(X,\OO_X)\}$ 给出。证明：这个函子是\trRepresentable{}的（被 $\Spec A$ 所\trRepresentation）。这就为 $\Spec A$ 的定义给出了另一条（得承认是有点怪的）动机，与 \S7.3.6 的那条密切相关。

\textbf{7.6.4.}\ \extrahard\ \textbf{群概形（或更一般地，\trCategory{}中的\trGroupObject{}）。}\quad（\S7.6 的其余部分，只该当作消遣来读。）我们又一次回到例~7.6.1。$X$ 上的函数可比一个集合要强：它们构成一个群。（的确，它们甚至构成一个\trRing，不过这件事我们日后再操心。）给定一个\trMorphism{} $X\to Y$，函数的\trPullback{} $\Gamma(Y,\OO_Y)\to\Gamma(X,\OO_X)$ 是一个\trGroupHom。于是我们应当指望 $\mathbb{A}^1$ 带有某种群式的结构。这就把我们引向\textbf{\trGroupScheme}（group scheme）的观念，或者更一般地，\trCategory{}中\textbf{\trGroupObject{}}（group object）的观念。我们现在就来定义它。

设 $\scr{C}$ 是一个带有\trFinalObject{} $Z$ 且带有\trProduct{}的\trCategory。（我们知道，$\mathit{Sch}$ 有\trFinalObject{} $Z=\Spec\mathbb{Z}$，见练习~7.3.J。至于它也有\trProduct，此事见 \S10.1，日后再说。不过在\mbox{练习~7.6.K} 中，我们将给出\trGroupObject{}的另一种刻画，它适用于任何\trCategory，所以此事不劳我们操心。）

$\scr{C}$ 中的一个\trGroupObject{}，是一个\trObject{} $X$ 连同三个\trMorphism：
\begin{itemize}[leftmargin=2.6em,itemsep=2pt,topsep=3pt]
  \item \emph{乘法}（multiplication）：$\mathrm{m}\colon X\times X\to X$；
  \item \emph{\trInverseElement}（inverse）：$\mathrm{i}\colon X\to X$；以及
  \item \emph{\trUnit}（identity element）：$\mathrm{e}\colon Z\to X$（\emph{不是恒等映射}）。
\end{itemize}
这些\trMorphism{}须满足以下几个条件。

（i）结合公理：
\[
\begin{tikzcd}[row sep=3.2em]
X\times X\times X \ar[r, "\mathrm{m}\times\mathrm{id}"] \ar[d, "\mathrm{id}\times\mathrm{m}"'] & X\times X \ar[d, "\mathrm{m}"] \\
X\times X \ar[r, "\mathrm{m}"'] & X
\end{tikzcd}
\]
交换。（这里的 $\mathrm{id}$ 指的是那个等同映射 $X\to X$。）

（ii）单位元公理：
\[
\begin{tikzcd}
X \ar[r, "\sim"] & Z\times X \ar[r, "\mathrm{e}\times\mathrm{id}"] & X\times X \ar[r, "\mathrm{m}"] & X
\end{tikzcd}
\]
与
\[
\begin{tikzcd}
X \ar[r, "\sim"] & X\times Z \ar[r, "\mathrm{id}\times\mathrm{e}"] & X\times X \ar[r, "\mathrm{m}"] & X
\end{tikzcd}
\]
都是恒等映射 $X\to X$。（这对应于群的那条公理：「乘以单位元便是恒等映射。」）

（iii）逆元公理：
\[
\begin{tikzcd}
X \ar[r, "\mathrm{i}\times\mathrm{id}"] & X\times X \ar[r, "\mathrm{m}"] & X
\end{tikzcd}
\]
与
\[
\begin{tikzcd}
X \ar[r, "\mathrm{id}\times\mathrm{i}"] & X\times X \ar[r, "\mathrm{m}"] & X
\end{tikzcd}
\]
都是经合成
\[
\begin{tikzcd}
X \ar[r] & Z \ar[r, "\mathrm{e}"] & X
\end{tikzcd}
\]
而得的那个映射。

作为一个动机，你可以验证：$\catSets$ 这个\trCategory{}中的\trGroupObject{}，其实与群是一回事。（这很能说明你是如何把某个观念「范畴化」的：你把它的诸公理按范畴论的路数写下来；倘若一切顺利，取 $\catSets$ 为特例时，你便回到了原来的观念。下面的练习会让你把这套功夫施之于「\trRing」的观念。）

\textbf{\trGroupScheme}（group scheme）定义为概形的\trCategory{}中的\trGroupObject{}。环 $A$ 上的群概形（或概形 $S$ 上的群概形）定义为 $A$-\trScheme（或 $S$-\trScheme）的\trCategory{}中的\trGroupObject{}。

\textbf{7.6.G.（练习）}\quad 请为 $\mathbb{A}^1_{\mathbb{Z}}$ 赋予群概形的结构：把三个\trStructureMorphism{}描述出来，并证明它们满足诸公理。（提示：这些态射不应令人意外。例如，\trInverseElement{}由 $t\mapsto-t$ 给出。注意：由\mbox{练习~7.6.E}(a) 我们已经知道积 $\mathbb{A}^1_{\mathbb{Z}}\times_{\mathbb{Z}}\mathbb{A}^1_{\mathbb{Z}}$ 存在。）

\textbf{7.6.H.（练习）}\quad 证明：若 $G$ 是\trCategory{} $\scr{C}$ 中的一个\trGroupObject{}，则对任何 $X\in\scr{C}$，$\Mor(X,G)$ 都带有群的结构，而且这个群结构在\trPullback{}下保持不变（也就是说，$\Mor(-,G)$ 是到群的\trCategory{} $\catGroups$ 的一个\trContravariantFunctor）。

\textbf{7.6.I.（练习）}\quad 证明：上一练习所描述的群结构，经由 \S7.6.1 的双射，把 $\mathbb{A}^1_{\mathbb{Z}}$ 上的群概形结构翻译为 $\Gamma(X,\OO_X)$ 上的群结构。

\textbf{7.6.J.（练习）}\quad 请定义\textbf{\trAbelian{}群概形}（abelian group scheme）与\textbf{\trRingScheme{}}（ring scheme）的观念。（毫无疑问，你会同时想明白如何在任何\trCategory{} $\scr{C}$ 中定义\trAbelian{}\trGroupObject{}与环对象的观念。读过 \S7.6.5 之后，你或许会发现一条处理这类问题的更省力的门径。）

\textbf{7.6.5.}\ \emph{群概形，更函子式地。}\quad 我们关于 \S7.6.1 中那个双射的「尊重群结构」性质的讨论，有一点令人不满意：我们先观察到右边（$X$ 上的函数）构成一个群，接着发展出群概形的公理，然后灵机一动琢磨出让 $\mathbb{A}^1_{\mathbb{Z}}$ 成为群概形的那些映射，最后才证明：这在双射的左边（$\Mor(X,\mathbb{A}^1)$）诱导出一个群结构，而这个群结构恰与右边（$X$ 上的函数）的群结构相对应。

这幅图景可以更干净地讲清楚，如下。概形\trValuedPoint{}的语言（定义~7.3.10）有这样一个好处：请注意，群概形的诸点本身未必构成一个群（想想 $\mathbb{A}^1_{\mathbb{Z}}$）。但\mbox{练习~7.6.H} 表明：群概形的诸 $Z$-\trValuedPoint（其中 $Z$ 是任何取定的\trScheme）确实构成一个群。

\textbf{7.6.K.（练习）}\quad 设我们有一个从 $\mathit{Sch}$（说实在的，任何\trCategory{}都行）到 $\catGroups$ 的\trContravariantFunctor{} $F$。再设 $F$ 与遗忘函子 $\catGroups\to\catSets$ 的复合被某个\trObject{} $Y$ 所\trRepresentation。证明：诸 $F(X)$ 上的群运算（当 $X$ 跑遍 $\mathit{Sch}$ 时）唯一地确定出满足群概形定义诸公理的 $\mathrm{m}\colon Y\times Y\to Y$、$\mathrm{i}\colon Y\to Y$、$\mathrm{e}\colon Z\to Y$，使得 $\Mor(X,Y)$ 上的群运算与 $F(X)$ 上的相同。

特别地，「群」的定义把「\trCategory{}中\trGroupObject{}」的定义逼到了我们手上。更一般地，你应当有这样的预期：任何一类可以被解读为「带有附加结构的集合」的对象，都该能装进这幅图景。

你应当把这个练习用到 $\mathbb{A}^1_{\mathbb{Z}}$ 上，看看你在\mbox{练习~7.6.G} 中找到的那些显式公式是怎样被逼出来的。

\textbf{7.6.L.（练习）}\quad 请把\mbox{练习~7.6.E} 的群概形中的诸映射 $\mathrm{m}$、$\mathrm{i}$ 与 $\mathrm{e}$ 具体算出来。

\textbf{7.6.M.（练习）}\quad 请解释：为什么一个\trCategory{}中诸\trGroupObject{}的\trProduct{}可以自然地解读为该\trCategory{}中的一个\trGroupObject{}。

\vkexsync{N}
\begin{exercise}\label{ex:morphism-of-group-schemes}
(a)~请定义\textbf{群概形的\trMorphism}（morphism of group schemes）。

(b)~回想：若 $A$ 是一个\trRing，则 $\mathrm{GL}_n(A)$（\textbf{$A$ 上的\trGeneralLinearGroup{}}（general linear group over $A$））是元素取自\trRing{} $A$ 的可逆 $n\times n$ 矩阵所成的群。请琢磨出\textbf{群概形 $\mathrm{GL}_n$（环 $A$ 上的）}的正确定义，并描述\textbf{\trDeterminant{}映射}（determinant map）$\det\colon\mathrm{GL}_n\to\mathbb{G}_m$。

(c)~请把下面这句话琢磨出道理来：「由 $t\mapsto t^n$ 给出的 $(\cdot^{n})\colon\mathbb{G}_m\to\mathbb{G}_m$ 是群概形的一个\trMorphism。」
\end{exercise}

\cref{ex:morphism-of-group-schemes}(a) 的语言暗示：群概形也构成一个\trCategory；你若乐意，大可去证明这件事。事实上，群概形的\trCategory{}还有一个\trZeroObject。它是什么？

\textbf{7.6.O.（练习：群概形态射的核）}\quad 设 $F\colon G_1\to G_2$ 是群概形的一个\trMorphism。考虑由 $X\mapsto\ker(\Mor(X,G_1)\to\Mor(X,G_2))$ 给出的\trContravariantFunctor{} $\mathit{Sch}\to\catGroups$。若这个函子是可表的——比如说被某个群概形 $G_0$ 所\trRepresentation——证明：$G_0\to G_1$ 是 $F$ 在群概形的\trCategory{}中的\trKernel。

\textbf{7.6.P.（练习）}\quad 证明：$(\cdot^{n})$ 的\trKernel{}（见\cref{ex:morphism-of-group-schemes}）是可表的。若 $n>0$，证明：在特征 $p$ 整除 $n$ 的域 $k$ 上，这个群概形是非\trReduced{}的。（这个群概形记作 $\mu_n$，很重要，尽管我们不会用到它。）

\textbf{7.6.Q.（练习）}\quad 证明：$\det\colon\mathrm{GL}_n\to\mathbb{G}_m$ 的\trKernel{}是可表的。这就是群概形 $\mathrm{SL}_n$。（你可以在 $\mathbb{Z}$ 上做这件事，也可以在域 $k$ 上做，甚至在任意\trRing{} $A$ 上做；其中的代数是一样的。）

\textbf{7.6.R.（练习）}\quad 请（用尽可能省事的办法）证明：$\mathbb{A}^1_k$ 是一个环 $k$-\trScheme。（这里的 $k$ 可以换成任何\trRing。）

\vkexsync{S}
\begin{exercise}\label{ex:group-scheme-action}
(a)~请定义（左）\textbf{群概形作用}（group scheme action）的观念，即一个群概形作用在一个概形上。

(b)~设 $A$ 是一个\trRing。证明：在 $A$ 上指定一个整值的\trGrading，等价于指定 $\mathbb{G}_m$ 在 $\Spec A$ 上的一个作用。（分次的这种解读出人意料地富有启发。警示：整值分次有两种可能的选取，二者各有各的道理；文献中两种都在使用。）
\end{exercise}

\textbf{7.6.6.}\ \emph{题外话：\trHopfAlgebra。}\quad 这里有一个我们不会用到的观念，不过现在把它定义出来毫不费事。设 $G=\Spec A$ 是一个\trAffine{}群概形，即一个作为\trAffineScheme{}的群概形。群概形的范畴论式定义，可以改用\trRing{} $A$ 的语言重新陈述。（这需要把注记~7.3.5 想透；见\mbox{练习~10.1.A}。）此时这些公理便定义出一个\textbf{\trHopfAlgebra{}}（Hopf algebra）。例如，我们有一条「\trComultiplication{}映射」（comultiplication map）$A\to A\otimes A$。

\textbf{7.6.T.（练习）}\quad 由于 $\mathbb{A}^1_{\mathbb{Z}}$ 是群概形，$\mathbb{Z}[t]$ 便带有一个\trHopfAlgebra{}结构。请描述\trComultiplication{}映射 $\mathbb{Z}[t]\to\mathbb{Z}[t]\otimes_{\mathbb{Z}}\mathbb{Z}[t]$。


\section{\extrahard\ \trGrassmannian{}：第一种构造}

\trGrassmannian{}是一种有用的几何构造，堪称「线性代数底下的那个几何对象」。在\trField{} $K=\mathbb{R}$ 或 $\mathbb{C}$ 上的（经典）几何中，正如\trProjectiveSpace{}参数化给定的 $n$ 维\trVectorSpace{}中的一维\trSubspace{}，\trGrassmannian{}参数化 $n$ 维空间中的 $k$ 维\trSubspace{}。\trGrassmannian{} $G(k,n)$ 是一个维数为 $k(n-k)$ 的\trManifold{}（在该\trField{}上）。\trManifold{}结构由如下方式给出。给定 $n$ 维空间的一组基 $(v_1,\dots,v_n)$，「大多数」$k$-平面（k-plane）都可以描述为由下面这 $k$ 个向量张成的\trSubspace{}：
\begin{equation*}
\left\langle v_1+\sum_{i=k+1}^{n}a_{1i}v_i,\;\;
v_2+\sum_{i=k+1}^{n}a_{2i}v_i,\;\;\dots,\;\;
v_k+\sum_{i=k+1}^{n}a_{ki}v_i\right\rangle.
\tag{7.7.0.1}\label{eq:grassmannian-basis-k-plane}
\end{equation*}
（你能描述哪些 $k$-平面\emph{不}属于这种形式吗？提示：\trRowReducedEchelonForm{}（row-reduced echelon form）。旁注：按\trNormalForm{}（normal form）对 $G(k,n)$ 作出的\trStratification{}（stratification），正是\trGrassmannian{}的\trSchubertCell{}（Schubert cells）分解。你或许能利用\trNormalForm{}证明：每个\trSchubertCell{}都\trIsomorphic{}于一个\trAffine{}空间。）这种形式的任何 $k$-平面都可以被唯一地如此描述。我们据此把这种形式的诸 $k$-平面等同于\trManifold{} $K^{k(n-k)}$（以诸 $a_{ji}$ 为\trCoordinate）。这是\trGrassmannian{}上的一大块\trAffine{}片（称作相对于这组基的「开\trSchubertCell{}」）。当诸 $v_i$ 变动时，这些片覆盖整个\trGrassmannian{}（为什么？），而且诸\trManifold{}结构彼此相合（这是一个更难的事实）。

我们现在于代数几何中、在一个\trRing{} $A$ 上来\emph{定义}\trGrassmannian。设 $v=(v_1,\dots,v_n)$ 是 $A^{\oplus n}$ 的一组基。更精确地说：诸 $v_i\in A^{\oplus n}$，并且由 $(a_1,\dots,a_n)\mapsto a_1v_1+\dots+a_nv_n$ 给出的映射 $A^{\oplus n}\to A^{\oplus n}$ 是一个\trIsomorphism。

\textbf{7.7.A.（练习）}\quad 证明：任何两组基都由 $A$ 上一个可逆的 $n\times n$ 矩阵相联系——所谓「可逆」，意即该矩阵元素取自 $A$，且其\trDeterminant{}是 $A$ 的一个可逆元。

对每一组这样的基 $v$，我们考虑\trScheme{} $U_v\cong\mathbb{A}^{k(n-k)}_A$，其\trCoordinate{}为诸 $a_{ji}$（$k+1\le i\le n$，$1\le j\le k$），我们把它想象成对应于由\eqref{eq:grassmannian-basis-k-plane}中诸向量张成的那一 $k$-平面。

\textbf{7.7.B.（练习）}\quad 给定两组基 $v$ 与 $w$，请说明如何沿适当的\trOpenSet{}把 $U_v$ 与 $U_w$ 粘合起来。你会发现下面这点或许带来方便：使用\trCoordinate{} $a_{ji}$（其 $i$ 从 $1$ 跑到 $n$，而不只是从 $k+1$ 跑到 $n$），但在 $i\le k$ 时强制约定 $a_{ji}=\delta_{ji}$（即 $i=j$ 时取 $1$，否则取 $0$）。这一约定与\trProjectiveSpace{}诸片上的\trCoordinate{} $x_{i/j}$ 相类似（\S4.4.9）。提示：$U_v$ 中相关的那个\trOpenSubset，将是某个\trDeterminant{}不消失之处。

\textbf{7.7.C.（练习/定义）}\quad 通过检验三重交，验证这些片（跑遍一切可能的基）能够粘合成单个\trScheme（\cref{ex:glue-arbitrary-schemes}）。这就是\trRing{} $A$ 上的\textbf{\trGrassmannian} $G(k,n)$。由于它可以解读为 $\mathbb{P}^{n-1}_A$ 中线性「$\mathbb{P}^{k-1}_A$」所成之空间，故也常写作 $G(k-1,n-1)$。（你将会看到这一记号确属明智，例如在\mbox{练习~12.2.N}中。）

尽管这个定义直白得令人愉快（\trGrassmannian{}被一片片 $\mathbb{A}^{k(n-k)}$ 覆盖是立见之事），或许也比我们在 \S4.4.9 中对\trProjectiveSpace{}的原始定义更「自然」（我们并不对基作选取；我们动用\emph{一切}基），它仍有数点不能令人满意。事实上，\trGrassmannian{}是\trProjective{}的；单看这个定义，这一点并不显然。此外，\trGrassmannian{}还自带一个到 $\mathbb{P}^{\binom{n}{k}-1}$ 的自然\trMorphism（它是\trClosedEmbedding{}的一个例子；这一观念我们将在定义~9.1.1 中遇到），称为\trPluckerEmbedding{}（Plücker embedding）。最后，$\mathrm{GL}_n$ 在「$n$ 维空间中的 $k$-平面」所成之空间上有一个作用，因此我们理应能在我们的代数化身中看到这个作用。我们将在 \S16.4 中处理这些问题，办法是给出一份更好的描述：把它作为\trModuliSpace{}。

\textbf{7.7.1.}\ \emph{（部分）\trFlagVariety{}。}\quad 正如\trGrassmannian{}「参数化」$n$ 维空间中的 $k$-平面，\textbf{\trFlagVariety{}}（flag variety）参数化「旗」（flags）：$n$ 维空间中\trSubspace{}的嵌套序列
\[
F_0\subset F_1\subset\dots\subset F_n,
\]
其中 $\dim F_i=i$。把这两者一并推广的，是与数据 $0\le a_1<\dots<a_\ell\le n$ 相关联的\textbf{\trPartialFlagVariety{}}（partial flag variety）这一观念，它参数化 $n$ 维空间中\trSubspace{}的嵌套序列
\[
F_{a_1}\subset\dots\subset F_{a_\ell},
\]
其中 $\dim F_{a_i}=a_i$。\S7.7 中的全部讨论，你都应当能够推广到这个场合。

